\chapter{Cap hitting estimates under critical-point conditioning}\label{sec:entrance}

In this section, we prove \cref{thm:ent-main}, the hitting estimate for the
caps under critical-point conditioning. It concerns the original
Langevin diffusion started from the uniform law $\nu_N$, in a field
conditioned to have a critical point of the center field $H_Q$ at a
prescribed point $u$, with prescribed value and radial derivative. The theorem
states that this diffusion reaches the set $\{x:|R(x,u)|\ge a\}$ before
time $e^{\gamma N}$ with averaged probability at most
$C(N^{1-d/2}+N^{3-d})+e^{-cN}$, for every finite mixture of minimum degree
$d\ge4$. \Cref{cor:ent-rectangle} makes the estimate uniform over a compact
rectangle of marks, with one time exponent. This corollary is the form of
the conditioned hitting estimate used in \cref{prop:cp-transfer}, where it is combined with the
determinant-weighted transfer of \cref{thm:kr-transfer} and the separation
criterion of \cref{prop:red-consumer}.

The section is self-contained. From earlier sections, it uses only the model
(\cref{def:pre-model}), the construction of the dynamics
(\cref{lem:pre-dynamics}), and the rotation covariance of the path kernels
(\cref{lem:red-kernels}). It does not use the static results of \cref{sec:thermal},
the counts of critical points and the concentration of the Gibbs measure on
bands in \cref{sec:purecover,sec:mixone,sec:mixpair,sec:mixbands}, or the Kac--Rice formula of
\cref{sec:kacrice}, and it needs no lower bound on the Hessian at the
critical point. The
spherical harmonic computations, the concentration estimate on the sphere,
and the comparison of resolvents that we need are proved here.

\paragraph{Standing notation.}
Throughout this section, $\xi(r)=\sum_{p\in\cP}c_pr^p$ is a finite mixture
with $c_p>0$ for $p\in\cP$, minimum degree $d=\min\cP\ge4$, and maximum
degree $\bar p=\max\cP$. We realize $H_N=\sum_{p\in\cP}H_{N,p}$, where the
$H_{N,p}$ are independent centered Gaussian homogeneous polynomials of degree
$p$ on $\R^N$ with $\E H_{N,p}(x)H_{N,p}(y)=Nc_pR(x,y)^p$. For fixed $N$,
the field is a function of the finite-dimensional Gaussian vector $J$ of the
coefficients of these polynomials. We write $H_J$ when this dependence
matters. The covariance identity $\E H_N(x)H_N(y)=N\xi(R(x,y))$, with
$R(x,y)=x\cdot y/N$, holds for all $x,y\in\R^N$. We fix $Q\in(0,1]$ and
$\beta\ge0$, and recall the center field and the radial derivative,
\begin{equation}\label{eq:ent-center}
 \begin{gathered}
 H_Q(u)=H_N(\sqrt Q\,u)=\sum_{p\in\cP}Q^{p/2}H_{N,p}(u),\\
 K_Q(u)=\frac{\dd}{\dd t}H_N(t\sqrt Q\,u)\Big|_{t=1}
       =\sum_{p\in\cP}pQ^{p/2}H_{N,p}(u),
 \end{gathered}
\end{equation}
and the parameters $\xi(Q)$, $a_Q=Q\xi'(Q)$, $b_Q=Q^2\xi''(Q)$, and
\[
 \Sigma_Q=\begin{pmatrix}\xi(Q)&a_Q\\ a_Q&a_Q+b_Q\end{pmatrix}.
\]
For $u\in S_N$ we write $G_u=\nabla_{S_N}H_Q(u)$, a vector in the tangent
space $u^\perp$. We write $\nabla_E$ for the Euclidean gradient,
$\nabla=\nabla_{S_N}$ for the gradient in the induced metric of $S_N$, and
$\Pi_x=I-xx^{\mathsf T}/N$ for the orthogonal projection onto $x^\perp$,
so that $\nabla H=\Pi_x\nabla_EH$ at $x\in S_N$. The diffusion with
Hamiltonian $H$ has generator $L_{\beta,H}=\Delta_{S_N}+\beta\nabla H\cdot
\nabla$ and invariant law $\pi^H$, and $\Pp^H_{\nu_N}$ denotes its path law
started from $\nu_N$. For $0<a<1$ and $u\in S_N$, we write
\[
 C_a(u)=\{x\in S_N:|R(x,u)|\ge a\}
\]
for the union of the two antipodal caps of latitude $a$ around $\pm u$, and
$I(a)=-\frac12\log(1-a^2)$.

We prescribe the values $H_Q(u)=Nh$ and $K_Q(u)=Nk$ and call the pair
$(h,k)$ of the energy and the radial derivative per spin the \emph{marks}. We fix bounds $\bar h,\bar k>0$ and call a pair $(h,k)$
\emph{admissible} if $|h|\le\bar h$ and $|k|\le\bar k$, and, when $\xi$ is
pure of degree $p$, also $k=ph$. Unless stated otherwise, $C$ and $c$ denote
positive constants that depend only on $\xi$, $Q$, $\beta$, $\bar h$, and
$\bar k$, and that may change from line to line. They never depend on $N$,
$u$, $(h,k)$, a point of $S_N$, or the realization of the field on the
regularity event of \cref{lem:ent-regularity}.

\section{The main estimate}\label{sec:ent-statement}

We first define the conditioned field. When $\xi$ is a non-pure mixture,
$\Sigma_Q$ is nonsingular (see \cref{prop:ent-regression}), and we define
polynomials $\theta_1,\theta_2,\theta_3$ by
\begin{equation}\label{eq:ent-theta}
 \bigl(\theta_1(r),\theta_2(r)\bigr)
 =\bigl(\xi(\sqrt Q\,r),\ \sqrt Q\,r\,\xi'(\sqrt Q\,r)\bigr)\Sigma_Q^{-1},
 \qquad
 \theta_3(r)=\frac{\sqrt Q\,\xi'(\sqrt Q\,r)}{a_Q}.
\end{equation}
When $\xi=c_pr^p$ is pure, we set instead
$\theta_1(r)=Q^{-p/2}r^p$, $\theta_2=0$, and
$\theta_3(r)=Q^{-p/2}r^{p-1}$; the last formula agrees with
\eqref{eq:ent-theta}.

\begin{definition}[The conditioned field]\label{def:ent-conditioned}
Let $u\in S_N$, and for $x\in\R^N$ let $r=x\cdot u/N$. The \emph{Gaussian
regression projection} of $H_N$ at $u$ is the random polynomial
\begin{equation}\label{eq:ent-deleted}
 D_u(x)=\theta_1(r)H_Q(u)+\theta_2(r)K_Q(u)+\theta_3(r)\,G_u\cdot x .
\end{equation}
For an admissible pair $(h,k)$, the \emph{conditioned field} is
\begin{equation}\label{eq:ent-conditioned}
 H^{h,k,u}=H_N+\delta_{h,k,u},\qquad
 \delta_{h,k,u}(x)=N\bigl(\theta_1(r)h+\theta_2(r)k\bigr)-D_u(x).
\end{equation}
We write $\E_{h,k,u}$ for the expectation under the law of $H^{h,k,u}$.
\end{definition}

We show in \cref{prop:ent-regression} below that the law of $H^{h,k,u}$ is
a version of the conditional law of $H_N$ given $H_Q(u)=Nh$, $K_Q(u)=Nk$,
and $G_u=0$. In the pure case, $K_Q(u)=pH_Q(u)$ identically, and the
conditioning is on $H_Q(u)=Nh$ and $G_u=0$. This is the version given by
Gaussian regression, and it depends continuously on $(h,k,u)$.

\begin{theorem}[Cap hitting bound from the uniform initial law]
\label{thm:ent-main}
Let $\xi$ be a finite mixture of minimum degree $d\ge4$, and fix
$Q\in(0,1]$, $\beta\ge0$, and $\bar h,\bar k>0$. Then there exists $a_*>0$
such that, for every $a\in(0,a_*]$, there exist $\gamma,c,C>0$ and $N_0$
with the following property. For every $N\ge N_0$, every $u\in S_N$, and
every admissible pair $(h,k)$,
\begin{equation}\label{eq:ent-main}
 \E_{h,k,u}\Bigl[\Pp^{H}_{\nu_N}\Bigl(\sup_{0\le t\le e^{\gamma N}}
   |R(X_t,u)|\ge a\Bigr)\Bigr]
 \le C\bigl(N^{1-d/2}+N^{3-d}\bigr)+e^{-cN}.
\end{equation}
The constant $a_*$ depends only on $\xi,Q,\beta,\bar h,\bar k$, and the
constants $\gamma,c,C,N_0$ depend only on these and on $a$.
\end{theorem}

In \eqref{eq:ent-main}, the field $H$ has the law of $H^{h,k,u}$, and $X$ is
the diffusion with generator $L_{\beta,H}$; we assume that its initial point
and its driving noise are independent of $H$. The event in
\eqref{eq:ent-main} is the event that the path enters $C_a(u)$ by time
$e^{\gamma N}$. Note that both
polynomial terms tend to zero, since $d\ge4$.

\begin{corollary}[Uniformity over a rectangle of marks]
\label{cor:ent-rectangle}
Let $\xi$, $Q$, and $\beta$ be as in \cref{thm:ent-main}, and let
$\mathcal K\subset\R^2$ be compact. Then there exists $a_*>0$ such that, for
every $a\in(0,a_*]$, there exist $\gamma,c,C>0$ and $N_0$, depending only on
$\xi,Q,\beta,\mathcal K,a$, such that for every $N\ge N_0$,
\begin{equation}\label{eq:ent-rectangle}
 \sup_{u\in S_N}\ \sup_{(h,k)\in\mathcal K}\
 \E_{h,k,u}\Bigl[\Pp^{H}_{\nu_N}\Bigl(\sup_{0\le t\le e^{\gamma N}}
   |R(X_t,u)|\ge a\Bigr)\Bigr]
 \le C\bigl(N^{1-d/2}+N^{3-d}\bigr)+e^{-cN}.
\end{equation}
When $\xi$ is pure of degree $p$, the supremum is over
$(h,k)\in\mathcal K$ with $k=ph$. If the admissible set of marks is empty,
the left side is taken to be zero.
\end{corollary}

\begin{proof}
Fix $\bar h,\bar k$ with $\mathcal K\subset[-\bar h,\bar h]\times
[-\bar k,\bar k]$ and apply \cref{thm:ent-main}. Its constants depend only
on the bounds $\bar h,\bar k$ on the marks, equivalently on the
bounds $|H_Q(u)|\le N\bar h$ and $|K_Q(u)|\le N\bar k$ on the conditioned
values, and on $\xi,Q,\beta,a$.
\end{proof}

\begin{remark}[Scope]\label{rem:ent-scope}
In \cref{prop:cp-transfer}, \cref{cor:ent-rectangle} is applied to a
compact rectangle of marks that contains the marks of every fixed
energy window $[m_N-M,m_N+M]$. Its constants, and in particular the time exponent $\gamma$,
do not depend on the half-width $M$ of the window chosen later, as long
as the marks of that window lie in $\mathcal K$. On the rectangle used there,
$k$ also stays at a fixed distance above $2\sqrt{b_Q}$, which the determinant moments of
\cref{sec:kacrice} require; the conditional law of $\nabla^2H_Q(u)$ under
$\E_{h,k,u}$ is computed in that section. The proof below does not use this
condition. We do not obtain uniformity as $Q\downarrow0$, as the coefficients of
$\xi$ degenerate, or as $a$ shrinks with $N$. For minimum degree $d=3$, the
term $N^{3-d}$ does not vanish, and our argument yields no hitting estimate.
\end{remark}

\section{The mechanism and the origin of the errors}
\label{sec:ent-mechanism}

We now explain why the dynamics started from $\nu_N$ cannot reach the caps
around a critical point of $H_Q$ with a large value within time
$e^{\gamma N}$, and where the two
polynomial errors in \eqref{eq:ent-main} come from. The discussion in this
section is informal; the precise statements and proofs follow in the next
sections.

\paragraph{The unconditioned field.}
We begin with the case where the field is not conditioned and the axis $u$
is uniform on $S_N$ and independent of $H_N$. A uniform point has overlap of
order $N^{-1/2}$ with $u$, and $\nu_N(C_a(u))\le e^{-NI(a)/2}$ for
large $N$. By rotation invariance, the Gibbs mass of $C_a(u)$, averaged over
the field and the axis, is also exponentially small. To turn this into a
bound that holds for the whole time interval $[0,e^{\gamma N}]$, we introduce
a resolvent. Let $\kappa\ge0$ be a killing rate supported near $C_a(u)$, and
let $\tau$ be the killing time, that is, the first time at which
$\int_0^t\kappa(X_v)\dd v$ exceeds an independent unit exponential variable.
We will see that $\E e^{-\lambda\tau}$ is
controlled by $\lambda^{-1}\int\kappa\dd\pi^H$, together with a cost for the
initial law. The relative entropy of $\nu_N$ with respect to $\pi^H$ is of
order $N$, and it enters only through a bound on the Fisher information of
the initial density after it is smoothed by the resolvent of the unperturbed
generator. With $\lambda=e^{-\gamma N}$, the term $\lambda^{-1}\int\kappa\dd\pi^H$
is of order
$e^{\gamma N-NI(a)/2}$, which is exponentially small when $\gamma$ is small.

\paragraph{The conditioning and the axis average.}
Conditioning $u$ to be a critical point of $H_Q$ with a large value changes
the whole field, and $H^{h,k,u}=H_N+\delta_{h,k,u}$. The correction is of
order $N$ near $u$, so we cannot treat it as a small perturbation of $H_N$.
The observation that makes a comparison possible is that the correction
vanishes to high order at the equator of $u$. The functions $\theta_1$ and
$\theta_2$ are of order $|r|^d$, and $\theta_3$ is of order $|r|^{d-1}$, so
that
\[
 |\delta_{h,k,u}(x)|\le CN|r|^d+C\sqrt N|r|^{d-1}|V_x(u)|,
 \qquad V_x(u)=N^{-1/2}\,G_u\cdot x,
\]
where $r=R(x,u)$. A fixed point $x$ has overlap of order $N^{-1/2}$ with a
uniform axis $u$, and $V_x(u)$ is of order one. This means that the
correction is of order $N^{1-d/2}$ for typical pairs $(x,u)$, and this is the
source of both polynomial errors. We therefore compare the conditioned
dynamics at a uniform axis with the unconditioned dynamics, using a smooth
truncation $\delta_\alpha$ of $\delta_{h,k,u}$ at a small fixed scale
$\alpha$ (see \eqref{eq:ent-cutoff} and the last paragraph of this section),
through the density $e^w$ with $w=\beta\delta_\alpha$. The truncation is
necessary: near $u$, the correction itself is of order $N$, and the axis
average of $e^{\beta\delta_{h,k,u}}$ can be exponentially large. For a
sufficiently small fixed $\alpha$, we will see that only two moments of
$e^w$ are needed, averaged over the axis for each fixed point $x$ and each
fixed field satisfying the regularity bounds of \cref{lem:ent-regularity}:
\[
 \eta=\E_u|e^w-1|=O(N^{1-d/2}),\qquad
 \epsilon=\E_u\bigl[e^w(1-e^{-w})^2\bigr]=O(N^{2-d}).
\]

\paragraph{The correlated gradient.}
This raises a difficulty, since the point $x$ in these averages may depend
on the field, and $G_u$ is a function of the same field. We cannot treat
$G_u$ as a Gaussian vector independent of $x$. To overcome this difficulty,
we fix the field and $x$ and average $u$ over a latitude
$\{u:R(x,u)=r\}$. We show in \cref{lem:ent-latitude} that $V_x$ concentrates
on each latitude, by the logarithmic Sobolev inequality on the sphere
(\cref{lem:ent-lsi}), and that its latitude mean can be computed exactly by
the Funk--Hecke formula. The top spherical harmonic component of each
$H_{N,p}$ contributes $\sqrt N|r|^{p-1}$ to this mean, and the lower
components contribute a bounded amount. Note that these lower components
have covariance weights $\omega_{p,j,N}=O(N^{-j})$, and hence pointwise
variances of order $N^{1-j}$ (\cref{lem:ent-harmonic}); they must be
kept.

\paragraph{The time horizon.}
The main obstacle is the length $e^{\gamma N}$ of the time interval. A
comparison of path laws, for example by Girsanov's formula, costs an amount
proportional to the time. To overcome this obstacle, we prove the
entropy--resolvent comparison of \cref{lem:ent-resolvent}. It compares the
Laplace transform $\E e^{-\lambda\tau}$ of the killing time for the
perturbed process with an integrated killing rate, in which the initial
density is smoothed by the resolvent of the unperturbed generator, at the cost
\begin{equation}\label{eq:ent-cost}
 \underbrace{\eta+\sqrt\epsilon}_{O(N^{1-d/2})}
 \quad+\quad
 \underbrace{\epsilon\,(1+\mathcal H_0)}_{O(N^{3-d})},
\end{equation}
where $\mathcal H_0=\int(1+f_0)\log(1+f_0)\dd\mu$, and $f_0$ is the density
of the initial law with respect to the reference measure $\mu$ of
\cref{lem:ent-resolvent}. In our application, $\mu$ is the joint law of the
field and of a point drawn from the Gibbs measure, $f_0\mu$ is the joint law
of the field and of the initial point, and $\mathcal H_0=O(N)$ by
\eqref{eq:ent-start-entropy}. Since $\int f_0\dd\mu=1$ and the function
$a\mapsto(1+a)\log(1+a)-a\log a$ is nonnegative and concave, Jensen's
inequality shows that $\mathcal H_0$ exceeds $\int f_0\log f_0\dd\mu$, the
entropy of $\nu_N$ relative to the Gibbs measure averaged over the field, by
at most $2\log2$. The point is that this cost is not multiplied by the
time: the factor $e^{\gamma N}$ appears only in the integrated killing rate,
which is exponentially small. At $d=4$, both terms in \eqref{eq:ent-cost} are
$O(N^{-1})$; at $d=3$, the second is of order one. In the proof of the
comparison, we test the resolvent equation of the killed diffusion against a weighted
function, and a commutator identity cancels a cross term that would
otherwise be of order one.

\paragraph{Truncation and the return to the true process.}
The comparison lemma requires a bounded perturbation whose size is small
compared with the entropy cost of the caps. For this reason, we truncate
$\delta_{h,k,u}$ by smooth cutoffs at scale $\alpha$ in both $r=R(x,u)$ and
$s=G_u\cdot x/N$. On the support of the cutoffs, the truncated correction is
at most $CN\alpha^d$, while a cap of latitude $\alpha/4$ costs
$NI(\alpha/4)\ge N\alpha^2/32$. Since $d>2$, we can choose $\alpha$ small and
then $\gamma$ small, so that the integrated killing rate remains exponentially
small. We then return to the true process using two facts. First, after a hit
of the cap $C_{\alpha/2}(u)$, the killing rate equals one during a time
interval of fixed length, except with exponentially small probability. Second, the true and
truncated processes coincide until $|r|$ or $|s|$ reaches $\alpha$. Reaching
$|r|=\alpha$ requires crossing $C_a(u)$ with $a=\alpha/2$ first. Reaching
$|s|=\alpha$ is exponentially unlikely. Indeed, $G_u$ is
exactly the Gaussian vector that the conditioning removes. It is independent
of the conditioned field and of the whole true trajectory, and $s_t$ is
Gaussian with variance at most $a_Q/N$ at each fixed time.

\paragraph{Organization.}
In \cref{sec:ent-regularity}, we compute the harmonic decomposition of the
field and prove regularity bounds with exponential tails.
\Cref{sec:ent-regression} constructs the conditioned field by regression and
bounds the correction. \Cref{sec:ent-multipliers} proves the latitude
concentration and the two moment bounds for $e^w$. \Cref{sec:ent-resolvent}
states and proves the entropy--resolvent comparison for abstract diffusions.
\Cref{sec:ent-truncated} applies it to the truncated process, and
\cref{sec:ent-boundary} shows that the true process does not reach the
cutoff boundary $\{|G_u\cdot x|/N=\alpha\}$. We assemble the
proof of \cref{thm:ent-main} in \cref{sec:ent-proof}.
\Cref{tab:ent-parameters} lists the parameters in the order in which they
are fixed; all of them are fixed before $N\to\infty$.

\begin{table}[ht]
\centering
\caption{Parameters in the proof of \cref{thm:ent-main}, in the order in
which they are chosen.}\label{tab:ent-parameters}
\begin{tabular}{@{}lll@{}}
\toprule
Parameter & Role & Depends on\\
\midrule
$\xi,Q,\beta,\bar h,\bar k$ & model, radius, temperature, mark bounds
  & given\\
$C_0$ & regularity constant (\cref{lem:ent-regularity}) & $\xi,Q$\\
$\alpha_1$ & largest cutoff scale (\cref{prop:ent-truncated})
  & $\xi,Q,\beta,\bar h,\bar k$\\
$a\in(0,\alpha_1/2]$, $\alpha=2a$ & latitude of the caps, cutoff scale
  & chosen\\
$\gamma$ & time exponent & the above and $a$\\
$C,c,N_0$ & constants in \eqref{eq:ent-main} & the above and $a$\\
\bottomrule
\end{tabular}
\end{table}

\section{Harmonic coefficients and regularity of the field}
\label{sec:ent-regularity}

The latitude mean of the tangent gradient term $G_u\cdot x$ involves every spherical harmonic
component of $H_{N,p}$, not only the top one. In this section, we compute
the variances of these components exactly and prove regularity bounds at
their correct scales.

Let $N\ge3$ and $\alpha_N=(N-2)/2$. The Gegenbauer polynomials
$C^{(\alpha_N)}_\ell$ are defined by
$(1-2rt+t^2)^{-\alpha_N}=\sum_{\ell\ge0}C^{(\alpha_N)}_\ell(r)t^\ell$.
Expanding the left side by the binomial series gives
\[
 C^{(\alpha_N)}_\ell(r)=\sum_{j=0}^{\lfloor\ell/2\rfloor}
 \frac{(-1)^j(\alpha_N)_{\ell-j}}{j!\,(\ell-2j)!}(2r)^{\ell-2j},
 \qquad C^{(\alpha_N)}_\ell(1)=\frac{(2\alpha_N)_\ell}{\ell!},
\]
where $(t)_j=t(t+1)\cdots(t+j-1)$ and $(t)_0=1$; see
\cite[Section~18.5(iii), (18.5.10)]{DLMF}. We define the normalized zonal
polynomials
\begin{equation}\label{eq:ent-zonal}
 \mathsf Z_{\ell,N}(r)=\frac{C^{(\alpha_N)}_\ell(r)}{C^{(\alpha_N)}_\ell(1)}
 =\frac{\ell!}{(2\alpha_N)_\ell}\sum_{j=0}^{\lfloor\ell/2\rfloor}
 \frac{(-1)^j(\alpha_N)_{\ell-j}}{j!\,(\ell-2j)!}(2r)^{\ell-2j},
\end{equation}
so that $\mathsf Z_{0,N}=1$ and $\mathsf Z_{\ell,N}(1)=1$. From the generating
function, we also obtain the differential equation
\begin{equation}\label{eq:ent-zonal-ode}
 (1-r^2)\mathsf Z''-(N-1)r\mathsf Z'+\ell(\ell+N-2)\mathsf Z=0,
 \qquad \mathsf Z=\mathsf Z_{\ell,N}.
\end{equation}

\begin{lemma}[Harmonic coefficients]\label{lem:ent-harmonic}
Let $N\ge3$.
\begin{enumerate}
\item For every integer $p\ge0$,
\begin{equation}\label{eq:ent-harmonic}
 r^p=\sum_{j=0}^{\lfloor p/2\rfloor}\omega_{p,j,N}\,
     \mathsf Z_{p-2j,N}(r),\qquad
 \omega_{p,j,N}=\frac{p!\,\Gamma(\alpha_N)\,(p-2j+\alpha_N)}
   {2^pj!\,\Gamma(p-j+\alpha_N+1)}\cdot
   \frac{(2\alpha_N)_{p-2j}}{(p-2j)!}.
\end{equation}
The coefficients $\omega_{p,j,N}$ are positive and sum to one.
\item For fixed $p$ and $0\le j\le p/2$, as $N\to\infty$,
\begin{equation}\label{eq:ent-harmonic-asymp}
 \omega_{p,j,N}=\frac{p!}{2^jj!\,(p-2j)!}N^{-j}\bigl(1+O_p(N^{-1})\bigr);
 \quad\text{in particular } \omega_{p,j,N}\le C_pN^{-j}.
\end{equation}
\item For fixed $\ell$ and $m$, the function $\mathsf Z_{\ell,N}(r)-r^\ell$
and its first $m$ derivatives are bounded by $C_{\ell,m}/N$ on $[-1,1]$. In
particular, for $\ell\ge1$ and $|r|\le1$,
\begin{equation}\label{eq:ent-zonal-derivative}
 |\mathsf Z'_{\ell,N}(r)|\le C_\ell\bigl(|r|^{\ell-1}+N^{-1}\bigr).
\end{equation}
\item Let $H_{N,p}|_{S_N}=\sum_{j}H_{p,p-2j}$ be the decomposition into
spherical harmonics, where $H_{p,\ell}$ has degree $\ell$. The fields
$H_{p,p-2j}$, for $p\in\cP$ and $0\le j\le p/2$, are independent centered
Gaussian fields on $S_N$, and
\begin{equation}\label{eq:ent-sector-cov}
 \E H_{p,p-2j}(x)H_{p,p-2j}(y)=Nc_p\,\omega_{p,j,N}\,
 \mathsf Z_{p-2j,N}(R(x,y)),\qquad x,y\in S_N.
\end{equation}
\end{enumerate}
\end{lemma}

For example, at $p=4$ we have
\[
 r^4=\frac{(N-1)(N+1)}{(N+2)(N+4)}\mathsf Z_{4,N}(r)
     +\frac{6(N-1)}{N(N+4)}\mathsf Z_{2,N}(r)+\frac{3}{N(N+2)}.
\]
In particular, the component of degree two of a quartic field has a
variance of order $N^{-1}$ relative to the top component.

\begin{proof}
We first recall the uniqueness of zonal harmonics. Fix $z\in S_N$. The
Laplacian $\Delta_{S_N}$ acts on functions of $r=R(x,z)$ by
$g\mapsto N^{-1}\bigl[(1-r^2)g''-(N-1)rg'\bigr]$. Then
\eqref{eq:ent-zonal-ode} implies that the function $x\mapsto\mathsf Z_{\ell,N}(R(x,z))$
is an eigenfunction of $\Delta_{S_N}$ with eigenvalue
$-\ell(\ell+N-2)/N$; since it is the restriction of a polynomial of degree
$\ell$, it is a spherical harmonic of degree $\ell$, and it is invariant
under the rotations fixing $z$. Conversely, a spherical harmonic of degree
$\ell$ invariant under these rotations is a polynomial in $R(x,z)$ of degree
at most $\ell$ that satisfies \eqref{eq:ent-zonal-ode}. The polynomial
solutions of \eqref{eq:ent-zonal-ode} of degree at most $\ell$ form a
one-dimensional space, since the coefficients are determined recursively by
the leading one. Such a harmonic is a scalar multiple of
$\mathsf Z_{\ell,N}(R(\cdot,z))$.

\paragraph{The expansion \eqref{eq:ent-harmonic}.}
By \eqref{eq:ent-zonal},
$\omega_{p,j,N}\mathsf Z_{p-2j,N}=p!\,\Gamma(\alpha_N)(p-2j+\alpha_N)
\bigl(2^pj!\,\Gamma(p-j+\alpha_N+1)\bigr)^{-1}C^{(\alpha_N)}_{p-2j}$.
Inserting the explicit sum for $C^{(\alpha_N)}_{p-2j}$, we find that the
coefficient of $r^{p-2s}$ in the right side of \eqref{eq:ent-harmonic} is
\[
 \frac{p!}{2^{2s}(p-2s)!}\sum_{j=0}^{s}
 \frac{(-1)^{s-j}(p-2j+\alpha_N)\,\Gamma(p-s-j+\alpha_N)}
 {j!\,(s-j)!\,\Gamma(p-j+\alpha_N+1)}.
\]
For $s=0$ this equals one. For $s\ge1$, recall the beta integral
\[
 \frac{\Gamma(p-s-j+\alpha_N)}{\Gamma(p-j+\alpha_N+1)}
 =\frac1{s!}\int_0^1t^{p-s-j+\alpha_N-1}(1-t)^s\dd t .
\]
Writing $y=1/t$, the binomial theorem and its derivative in $y$ give
\[
 \sum_{j=0}^s\frac{(-1)^{s-j}(p+\alpha_N-2j)\,t^{-j}}{j!\,(s-j)!}
 =\frac{t^{-s}(1-t)^{s-1}}{s!}\bigl[(p+\alpha_N)(1-t)-2s\bigr].
\]
Inserting these two identities, we find that the sum over $j$ equals
$(s!)^{-2}$ times the integral over $[0,1]$ of
\[
 t^{p-2s+\alpha_N-1}(1-t)^{2s-1}\bigl[(p+\alpha_N)(1-t)-2s\bigr]
 =\frac{\dd}{\dd t}\Bigl[t^{p-2s+\alpha_N}(1-t)^{2s}\Bigr].
\]
Both boundary values vanish, because $p-2s+\alpha_N>0$ and $2s>0$. This
proves \eqref{eq:ent-harmonic}. The coefficients are positive by their
formula, and evaluating \eqref{eq:ent-harmonic} at $r=1$ shows that they sum
to one.

\paragraph{Asymptotics.}
For fixed $p$ and $j$, each gamma ratio and Pochhammer symbol in
$\omega_{p,j,N}$ is a finite product of factors $\alpha_N+c$, or of factors
$2\alpha_N+c$, with fixed constants $c$, or the reciprocal of such a
product. The substitution $\alpha_N+c=\alpha_N(1+c/\alpha_N)$ gives
\[
 \omega_{p,j,N}=\frac{p!}{2^{2j}j!\,(p-2j)!}\,\alpha_N^{-j}
 \bigl(1+O_p(\alpha_N^{-1})\bigr),
\]
which is \eqref{eq:ent-harmonic-asymp}, since $\alpha_N^{-j}=2^jN^{-j}
(1+O_j(N^{-1}))$. In the same way, the leading coefficient of
$\mathsf Z_{\ell,N}$ is $2^\ell(\alpha_N)_\ell/(2\alpha_N)_\ell=1+O_\ell
(N^{-1})$, and its coefficient of $r^{\ell-2j}$ is $O_\ell(N^{-j})$ for
$j\ge1$. Differentiating this polynomial of fixed degree yields (3).

\paragraph{The harmonic components.}
The decomposition of $H_{N,p}|_{S_N}$ into spherical harmonics is given by
linear projections, and the components are jointly Gaussian with the
original fields. Their covariances are obtained by applying the projections
in both variables of the kernel $Nc_pR(x,y)^p
=Nc_p\sum_j\omega_{p,j,N}\mathsf Z_{p-2j,N}(R(x,y))$. With one variable
fixed, $\mathsf Z_{\ell,N}(R(x,y))$ is a spherical harmonic of degree $\ell$
in the other variable. Harmonics of different degrees are orthogonal in
$L^2(\nu_N)$, as they are eigenfunctions of $\Delta_{S_N}$ with the distinct
eigenvalues $-\ell(\ell+N-2)/N$. The projection onto degree $\ell'$ in both
variables keeps the term with $\ell=\ell'$ and removes the others. This
proves \eqref{eq:ent-sector-cov}, and it shows that components of different
degrees, or of different $p$, are uncorrelated. Being jointly Gaussian, they
are independent.
\end{proof}

The next lemma is a refinement of \cref{lem:pre-field} for the finite
mixtures of this section: it also controls the lower harmonic components, at
their own scale. For a component $H_{p,\ell}$, we work with its harmonic
homogeneous extension of degree $\ell$ to $\R^N$.

\begin{lemma}[Regularity with exponential tails]\label{lem:ent-regularity}
There exist $C_0$ and $N_0$, depending only on $\xi$ and $Q$, and events
$\Omega_N$ determined by the coefficient vector $J$, such that
$\Pp(\Omega_N^c)\le e^{-N}$ for $N\ge N_0$, each $\Omega_N$ is invariant under
the rotations $H\mapsto H\circ g^{-1}$, $g\in O(N)$, and the following holds
on $\Omega_N$.
\begin{enumerate}
\item For every $p\in\cP$ and $m\in\{0,1,2\}$,
$\sup_{x\in S_N}\norm{\nabla_E^mH_{N,p}(x)}_{\rm op}\le C_0N^{1-m/2}$.
\item For every $p\in\cP$ and $0\le j\le p/2$,
$\sup_{x\in S_N}|H_{p,p-2j}(x)|\le C_0N^{1-j/2}$.
\item With $\norm{\cdot}_\infty$ denoting the supremum over $S_N$,
\begin{equation}\label{eq:ent-regularity}
 \begin{gathered}
 \norm{H_N}_\infty+\norm{H_Q}_\infty+\norm{K_Q}_\infty\le C_0N,\\
 \norm{\nabla_EH_N}_\infty+\norm{\nabla_EH_Q}_\infty
   +\norm{\nabla H_N}_\infty+\norm{\nabla H_Q}_\infty\le C_0\sqrt N,\\
 \norm{\nabla_E^2H_N}_\infty+\norm{\nabla_E^2H_Q}_\infty
   +\norm{\nabla^2H_N}_\infty+\norm{\nabla^2H_Q}_\infty\le C_0,
 \end{gathered}
\end{equation}
where the norms of second derivatives are operator norms.
\end{enumerate}
\end{lemma}

\begin{proof}
We first treat a component of degree zero, which occurs as $H_{p,0}$ for
even $p$. It is a centered Gaussian constant $Z$ with $\Var Z=N\sigma^2$,
where $\sigma^2=c_p\omega_{p,p/2,N}$. The Gaussian tail bound gives
$\Pp(|Z|>2L\sigma N)\le2e^{-2L^2N}$ for $L>0$, and on the complementary
event $|Z|\le2L\sigma N\le CLN^{1-p/4}$, which is the value bound required
in (2) with $j=p/2$; all derivatives of positive order vanish.

Now let $T$ be an isotropic centered Gaussian homogeneous polynomial of
degree $1\le\ell\le\bar p$ on $\R^N$, with $\E T(x)^2=N\sigma^2$ for $|x|^2=N$. By
homogeneity and isotropy, $\E T(y)^2=N\sigma^2(|y|^2/N)^\ell$ for all $y$.
Let $\mathcal T$ be the symmetric $\ell$-linear form with
$T(y)=\mathcal T(y,\dots,y)$. For unit vectors $v_1,\dots,v_\ell$, the
polarization identity
\[
 \mathcal T(v_1,\dots,v_\ell)=\frac1{2^\ell\ell!}
 \sum_{\varepsilon\in\{\pm1\}^\ell}\Bigl(\prod_i\varepsilon_i\Bigr)
 T\Bigl(\sum_i\varepsilon_iv_i\Bigr)
\]
and $|\sum_i\varepsilon_iv_i|\le\ell$ give
$\E\mathcal T(v_1,\dots,v_\ell)^2\le D_\ell^2\sigma^2N^{1-\ell}$, where
$D_\ell=\ell^\ell/\ell!$. Let $\mathcal N$ be a $(4\bar p)^{-1}$-net of the
unit sphere; comparing disjoint balls shows that we may take
$|\mathcal N|\le(1+8\bar p)^N$. Replacing the arguments one at a time by net
points, and using that each replacement costs at most
$(4\bar p)^{-1}\norm{\mathcal T}_{\rm op}$, we find
$\norm{\mathcal T}_{\rm op}\le\frac43\max_{v_i\in\mathcal N}
|\mathcal T(v_1,\dots,v_\ell)|$. The Gaussian tail bound and a union bound
over the $(1+8\bar p)^{\ell N}$ tuples then give, for $L>0$,
\begin{equation}\label{eq:ent-net}
 \Pp\bigl(\norm{\mathcal T}_{\rm op}>2LD_\ell\sigma N^{1-\ell/2}\bigr)
 \le2\exp\bigl(-N\bigl[L^2/2-\bar p\log(1+8\bar p)\bigr]\bigr).
\end{equation}
For $|x|=\sqrt N$ and $0\le m\le\ell$, the $m$th derivative of $T$ at $x$
is $(\ell)_m^{\downarrow}\,\mathcal T(x,\dots,x,\cdot,\dots,\cdot)$, where
$(\ell)_m^{\downarrow}=\ell(\ell-1)\cdots(\ell-m+1)$, and has operator norm
at most $(\ell)_m^{\downarrow}\norm{\mathcal T}_{\rm op}N^{(\ell-m)/2}$. On
the complement of the event in \eqref{eq:ent-net}, it is at most
$2LD_\ell(\ell)_m^{\downarrow}\sigma N^{1-m/2}$. Derivatives of order
$m>\ell$ vanish.

We use these bounds for each $H_{N,p}$, with $\sigma^2=c_p$, and for the
harmonic extension of each $H_{p,p-2j}$, with $\sigma^2=c_p\omega_{p,j,N}\le CN^{-j}$
by \eqref{eq:ent-sector-cov} and \eqref{eq:ent-harmonic-asymp}. This is a
fixed finite number $M_0$ of polynomials, those of degree zero included.
We choose $L$ with
$L^2/2\ge\bar p\log(1+8\bar p)+\log(2M_0)+2$. The union of the $M_0$
exceptional events then has probability at most $e^{-N}$ for large $N$; any
fixed exponential rate is possible by enlarging $L$. We define $\Omega_N$ by
the resulting bounds on the suprema in (1) and (2), with a suitable $C_0$.
These suprema are continuous functions of $J$, which makes $\Omega_N$
Borel, and they are invariant under rotations, since the decomposition into components
commutes with rotations.

Part (3) follows by summing the components with the weights $1$,
$Q^{p/2}$, and $pQ^{p/2}$ of \eqref{eq:ent-center}. For the spherical
derivatives, recall that $\nabla T=\Pi_x\nabla_ET$ and
\[
 \nabla^2T(x)=\Pi_x\nabla_E^2T(x)\Pi_x
 -\frac{x\cdot\nabla_ET(x)}{N}\,I_{x^\perp},
\]
and that $x\cdot\nabla_ET=\ell T$ for a homogeneous polynomial of
degree $\ell$. Increasing $C_0$ once completes the proof.
\end{proof}

\section{Exact regression of the whole field}\label{sec:ent-regression}

In this section, we show that the field of \cref{def:ent-conditioned} has
the conditional law, and we bound the correction $\delta_{h,k,u}$. The
regression is performed on the whole field $H_N$, although the conditioned
quantities concern the center field at radius $\sqrt Q$; the dynamics are
always driven by the field on $S_N$ itself.

\begin{proposition}[Whole-field regression]\label{prop:ent-regression}
Let $u\in S_N$, and let $r=x\cdot u/N$ for $x\in\R^N$.
\begin{enumerate}
\item If $\xi$ is a non-pure mixture, then $\det\Sigma_Q>0$, and the
coefficient vector of the random polynomial $H_N-D_u$ is independent of the
vector $(H_Q(u),K_Q(u),G_u)$. If $\xi$ is pure of degree $p$, then
$K_Q(u)=pH_Q(u)$, and the coefficient vector of $H_N-D_u$ is independent of
$(H_Q(u),G_u)$. In both cases, the law of $H^{h,k,u}$ is a version of the
conditional law of $H_N$ given $H_Q(u)=Nh$, $K_Q(u)=Nk$, and $G_u=0$.
\item For $|r|\le1$,
\begin{equation}\label{eq:ent-theta-orders}
 \begin{gathered}
 |\theta_1(r)|+|\theta_2(r)|\le C|r|^d,\qquad
 |\theta_3'(r)|\le C|r|^{d-2},\\
 |\theta_1'(r)|+|\theta_2'(r)|+|\theta_3(r)|\le C|r|^{d-1}.
 \end{gathered}
\end{equation}
\item The correction has spherical gradient
\begin{equation}\label{eq:ent-correction-gradient}
 \begin{split}
 \nabla\delta_{h,k,u}(x)={}&\Bigl[\theta_1'(r)\bigl(Nh-H_Q(u)\bigr)
 +\theta_2'(r)\bigl(Nk-K_Q(u)\bigr)-\theta_3'(r)\,G_u\cdot x\Bigr]
 \frac{\Pi_xu}{N}\\
 &-\theta_3(r)\Pi_xG_u
 \end{split}
\end{equation}
at $x\in S_N$. On $\Omega_N$, for every admissible $(h,k)$ and $x\in S_N$,
\begin{equation}\label{eq:ent-delta-bounds}
 |\delta_{h,k,u}(x)|\le CN|r|^d+C|r|^{d-1}|G_u\cdot x|,
 \qquad
 \norm{\nabla\delta_{h,k,u}}_\infty\le C\sqrt N .
\end{equation}
\end{enumerate}
\end{proposition}

\begin{proof}
We begin with the covariances. Recall that
$\E H_N(x)H_N(y)=N\xi(x\cdot y/N)$ for all $x,y\in\R^N$. Differentiating in
$y$ along $t\sqrt Q\,u$ and along unit vectors $e\perp u$, we find, for
every $x\in\R^N$,
\begin{equation}\label{eq:ent-covariances}
 \begin{gathered}
 \Cov\bigl((H_Q(u),K_Q(u))\bigr)=N\Sigma_Q,\qquad
 \Cov(G_u)=a_QI_{u^\perp},\\
 \Cov\bigl((H_Q(u),K_Q(u)),G_u\bigr)=0,\qquad
 \Cov\bigl(H_N(x),G_u\bigr)=\sqrt Q\,\xi'(\sqrt Q\,r)\Pi_ux,\\
 \Cov\bigl(H_N(x),(H_Q(u),K_Q(u))\bigr)
   =N\bigl(\xi(\sqrt Q\,r),\sqrt Q\,r\,\xi'(\sqrt Q\,r)\bigr).
 \end{gathered}
\end{equation}
For example, the identity for $\Cov(H_N(x),G_u)$ follows by
differentiating $N\xi(\sqrt Q\,x\cdot v/N)$ in $v$ at $v=u$ in a direction
$e\perp u$, which yields $\sqrt Q\,\xi'(\sqrt Q\,r)\,x\cdot e$. The covariance of $G_u$ follows
from $\partial_e\partial_{e'}N\xi(Qv\cdot v'/N)=Q\xi'(Q)\,e\cdot e'$ at
$v=v'=u$, since $u\cdot e=u\cdot e'=0$. The same orthogonality shows that
$(H_Q(u),K_Q(u))$ and $G_u$ are uncorrelated. With $w_p=c_pQ^p$, we have
$\xi(Q)=\sum_pw_p$, $a_Q=\sum_ppw_p$, and $a_Q+b_Q=\sum_pp^2w_p$, so that
\[
 \det\Sigma_Q=\Bigl(\sum_pw_p\Bigr)\Bigl(\sum_pp^2w_p\Bigr)
 -\Bigl(\sum_ppw_p\Bigr)^2>0
\]
by the Cauchy--Schwarz inequality, with equality exactly when only one degree
occurs.

Consider a non-pure mixture. By \eqref{eq:ent-covariances}, $D_u(x)$ is the
orthogonal projection of $H_N(x)$ onto the span of the coordinates of
$(H_Q(u),K_Q(u),G_u)$. Indeed, the coefficients of $H_Q(u)$ and $K_Q(u)$
are $N\bigl(\xi(\sqrt Q\,r),\sqrt Q\,r\,\xi'(\sqrt Q\,r)\bigr)(N\Sigma_Q)^{-1}
=(\theta_1(r),\theta_2(r))$, the coefficient vector of $G_u$ is
$\sqrt Q\,\xi'(\sqrt Q\,r)\Pi_ux/a_Q$, and $\Pi_ux\cdot G_u=x\cdot
G_u$. The family $\{(H_N-D_u)(x):x\in\R^N\}$ and the vector
$(H_Q(u),K_Q(u),G_u)$ are jointly Gaussian with zero cross-covariances, and
they are therefore independent. A polynomial of degree at most $\bar p$ is
determined by its values at finitely many points, and its coefficient vector
is a linear function of these values. This proves the independence in
(1). Since
$H_N=(H_N-D_u)+D_u$ and $D_u$ is linear in the conditioning vector, replacing
$(H_Q(u),K_Q(u),G_u)$ by $(Nh,Nk,0)$ in $D_u$ produces a version of the
conditional law; this is the field \eqref{eq:ent-conditioned}. In the pure
case, $H_N=H_{N,p}$ is homogeneous of degree $p$, which implies that
$K_Q(u)=pH_Q(u)$. We have $\Var H_Q(u)=Nc_pQ^p$ and
$\Cov(H_N(x),H_Q(u))=Nc_pQ^{p/2}r^p$. By \eqref{eq:ent-covariances}, the
regression coefficients on $H_Q(u)$ and $G_u$ are $Q^{-p/2}r^p$ and
$Q^{-p/2}r^{p-1}$. The same argument applies, and no singular matrix is
inverted.

For (2), note that $\theta_1,\theta_2$ are linear combinations of
$\xi(\sqrt Q\,r)$ and $\sqrt Q\,r\,\xi'(\sqrt Q\,r)$, whose lowest power of
$r$ is $r^d$, and that $\theta_3$ is a multiple of $\xi'(\sqrt Q\,r)$, whose
lowest power is $r^{d-1}$. These are fixed polynomials of degree at most
$\bar p$, which implies \eqref{eq:ent-theta-orders}.

For (3), the identity \eqref{eq:ent-correction-gradient} is the product rule,
with $\nabla r=\Pi_xu/N$ and $\nabla(G_u\cdot x)=\Pi_xG_u$. On
$\Omega_N$, we have $|Nh-H_Q(u)|\le(\bar h+C_0)N$, $|Nk-K_Q(u)|\le
(\bar k+C_0)N$, $|G_u|\le C_0\sqrt N$, and $|\Pi_xu|/N\le N^{-1/2}$.
Inserting these bounds and \eqref{eq:ent-theta-orders} into
\eqref{eq:ent-conditioned} and \eqref{eq:ent-correction-gradient}, and using
$|G_u\cdot x|\le C_0N$, we obtain \eqref{eq:ent-delta-bounds}. These bounds
control the conditional mean $N(\theta_1h+\theta_2k)$, which is of order $N$
near $u$, as well as the random part of the correction.
\end{proof}

\section{Latitude concentration and moments of the density ratio}
\label{sec:ent-multipliers}

In this section, we bound two moments of the density $e^w$ of the
truncated perturbation, averaged over a uniform axis, for every fixed field
on $\Omega_N$ and every fixed point $x$. We begin with the concentration
inequality on the sphere. This is the curvature criterion of Bakry and \'Emery
\cite[Corollary~2, p.~199; Example~(b), p.~200]{BakryEmery1985} for the
logarithmic Sobolev inequality, followed by the Herbst argument; see also
\cite[Sections~5.4 and~5.7]{BakryGentilLedoux2014}.
The optimal constant on the sphere is due to Mueller and Weissler
\cite{MuellerWeissler1982}; see also \cite[p.~200]{BakryEmery1985}. The
curvature constant used below suffices for our purposes.

\begin{lemma}[Concentration on the sphere]\label{lem:ent-lsi}
Let $m\ge2$, let $\sigma_m$ be the normalized surface measure on the unit
sphere $S^m\subset\R^{m+1}$, and let $g:S^m\to\R$ be $L$-Lipschitz for the
geodesic distance. Then
\begin{equation}\label{eq:ent-subgaussian}
 \int e^{t(g-\int g\dd\sigma_m)}\dd\sigma_m
 \le\exp\Bigl(\frac{L^2t^2}{2(m-1)}\Bigr),\qquad t\in\R .
\end{equation}
\end{lemma}

\begin{proof}
For a function $q>0$, let $\operatorname{Ent}(q)=\int q\log q\dd\sigma_m
-\int q\dd\sigma_m\log\int q\dd\sigma_m$ and
$\mathcal I(q)=\int|\nabla q|^2/q\dd\sigma_m$. We first show that
\begin{equation}\label{eq:ent-lsi}
 \operatorname{Ent}(q)\le\frac{\mathcal I(q)}{2(m-1)}
\end{equation}
for smooth $q>0$. Let $q_t=e^{t\Delta}q$ be the heat flow on $S^m$, which
stays smooth and positive. By integration by parts,
we have $\frac{\dd}{\dd t}\operatorname{Ent}(q_t)=-\mathcal I(q_t)$. Writing
$v=\log q_t$, we have $\partial_tv=\Delta v+|\nabla v|^2$ and
$\mathcal I(q_t)=\int q_t|\nabla v|^2$. Differentiating, we find
\[
 \frac{\dd}{\dd t}\mathcal I(q_t)
 =\int q_t\Delta|\nabla v|^2+2\int q_t\nabla v\cdot\nabla\Delta v
  +2\int q_t\nabla v\cdot\nabla|\nabla v|^2 ,
\]
where we moved the Laplacian in the first term. The divergence theorem for
$q_t\nabla|\nabla v|^2$, with $\nabla q_t=q_t\nabla v$, shows that the third
term equals $-2\int q_t\Delta|\nabla v|^2$. By the Bochner identity
$\frac12\Delta|\nabla v|^2=\nabla v\cdot\nabla\Delta v+|\nabla^2v|^2
+\operatorname{Ric}(\nabla v,\nabla v)$ and $\operatorname{Ric}=(m-1)g_{S^m}$,
we conclude that
\[
 \frac{\dd}{\dd t}\mathcal I(q_t)
 =-2\int q_t\bigl(|\nabla^2v|^2+(m-1)|\nabla v|^2\bigr)
 \le-2(m-1)\mathcal I(q_t).
\]
Hence $\mathcal I(q_t)\le e^{-2(m-1)t}\mathcal I(q)$. By the spherical harmonic
expansion, $q_t$ converges uniformly to $\int q\dd\sigma_m$, and
$\operatorname{Ent}(q_t)\to0$. Integrating in $t$ gives
\eqref{eq:ent-lsi}.

Now let $g$ be smooth with $|\nabla g|\le L$, and let
$\psi(t)=\log\int e^{tg}\dd\sigma_m$. Applying \eqref{eq:ent-lsi} to
$q=e^{tg}$, for which $\mathcal I(q)\le L^2t^2\int e^{tg}$, we find
$t\psi'(t)-\psi(t)\le L^2t^2/(2(m-1))$. The derivative of $\psi(t)/t$ is
at most $L^2/(2(m-1))$, and $\psi(t)/t\to\int g\dd\sigma_m$ as
$t\downarrow0$. Integrating, we obtain \eqref{eq:ent-subgaussian} for $t>0$;
the case $t<0$ follows by applying this to $-g$. A Lipschitz function is a uniform
limit of smooth functions with the same Lipschitz constant, and dominated
convergence completes the proof.
\end{proof}

We now fix $x\in S_N$ and let $u$ be uniform on $S_N$. Then $r=R(x,u)$ has
density
\begin{equation}\label{eq:ent-axis-density}
 p_N(r)=c_N(1-r^2)^{(N-3)/2},\quad -1<r<1,\qquad
 c_N=\frac{\Gamma(N/2)}{\sqrt\pi\,\Gamma((N-1)/2)}\le\sqrt N,
\end{equation}
by spherical coordinates and the beta integral; the bound on $c_N$ follows
from $\Gamma(z+\frac12)\le\sqrt z\,\Gamma(z)$. Conditionally on $r$, we can
write $u=rx+\sqrt{1-r^2}\,y$, where $y$ is uniform on the latitude sphere
$S_N\cap x^\perp$, which has dimension $N-2$ and radius $\sqrt N$. We write
$\E_{\rm lat}$ for the expectation over $y$ at fixed $r$, and define
\begin{equation}\label{eq:ent-V}
 V_x(u)=N^{-1/2}\,G_u\cdot x=N^{-1/2}\,\nabla H_Q(u)\cdot x .
\end{equation}

\begin{lemma}[The tangent gradient on a latitude]\label{lem:ent-latitude}
Let $N\ge6$, let the field lie in $\Omega_N$, and fix $x\in S_N$. For
$-1<r<1$, let $m_x(r)=\E_{\rm lat}V_x(u)$ and $Z=V_x(u)-m_x(r)$. Then
\begin{equation}\label{eq:ent-latitude}
 |m_x(r)|\le C\bigl(\sqrt N|r|^{d-1}+1\bigr),\qquad
 \E_{\rm lat}e^{tZ}\le e^{Ct^2},\quad t\in\R .
\end{equation}
\end{lemma}

The point $x$ in this lemma may be any function of the field: the lemma
fixes the field and $x$ before averaging over the axis.

\begin{proof}
We first prove the concentration. Along a unit-speed curve in $S_N$, the
derivative of $\Pi_u\nabla_EH_Q(u)$ has norm at most $2|\nabla_EH_Q(u)|/\sqrt N+\norm{\nabla^2_EH_Q(u)}_{\rm op}
\le3C_0$ by \eqref{eq:ent-regularity}. Because $|x|=\sqrt N$, the function
$V_x$ is $3C_0$-Lipschitz on $S_N$. The map $y\mapsto rx+\sqrt{1-r^2}\,y$ is
$1$-Lipschitz. After rescaling the latitude sphere to the unit sphere
$S^{N-2}$, the function $y\mapsto V_x(u)$ is $3C_0\sqrt
N$-Lipschitz. Then \cref{lem:ent-lsi} with $m=N-2$ implies that
$\E_{\rm lat}e^{tZ}\le\exp(9C_0^2Nt^2/(2(N-3)))\le e^{Ct^2}$.

We next compute the mean, keeping all harmonic components. Let $Y$ be a
spherical harmonic of degree $\ell$ on $S_N$. Its average over the rotations
fixing $x$ is a spherical harmonic of degree $\ell$ invariant under these
rotations, with the value $Y(x)$ at $x$. By the uniqueness of zonal
harmonics recalled in the proof of \cref{lem:ent-harmonic}, it equals
$Y(x)\mathsf Z_{\ell,N}(R(x,\cdot))$. Evaluated on the latitude of overlap
$r$, this gives $\E_{\rm lat}Y(u)=\mathsf Z_{\ell,N}(r)\,Y(x)$, the
latitude-averaging form of the Funk--Hecke identity. For the uniqueness of
zonal harmonics used in this argument, see also
\cite[Theorem~2.8, pp.~19--20]{AtkinsonHan2012}.
Using \eqref{eq:ent-center} and \cref{lem:ent-harmonic}(4), we obtain
\begin{equation}\label{eq:ent-funk-hecke}
 \E_{\rm lat}H_Q(u)=\sum_{p\in\cP}Q^{p/2}\sum_{j=0}^{\lfloor p/2\rfloor}
 \mathsf Z_{p-2j,N}(r)\,H_{p,p-2j}(x).
\end{equation}
For fixed $y$, we have $\partial_ru=x-ry/\sqrt{1-r^2}=(x-ru)/(1-r^2)$, and
$(x-ru)\cdot u=0$. Since $G_u$ is tangent at $u$, we have
\[
 \frac{\dd}{\dd r}H_Q(u)=\nabla_EH_Q(u)\cdot\frac{x-ru}{1-r^2}
 =\frac{G_u\cdot x}{1-r^2}=\frac{\sqrt N\,V_x(u)}{1-r^2}.
\]
Since the law of $y$ does not depend on $r$, we may differentiate
\eqref{eq:ent-funk-hecke} under the expectation. We find
\begin{equation}\label{eq:ent-latitude-mean}
 m_x(r)=\frac{1-r^2}{\sqrt N}\sum_{p\in\cP}Q^{p/2}
 \sum_{j=0}^{\lfloor p/2\rfloor}\mathsf Z'_{p-2j,N}(r)\,H_{p,p-2j}(x).
\end{equation}
For $j=0$, \cref{lem:ent-regularity}(2) and \eqref{eq:ent-zonal-derivative}
bound the term by $C\sqrt N|r|^{p-1}+CN^{-1/2}$. For $j\ge1$, the component
satisfies $|H_{p,p-2j}(x)|\le C_0\sqrt N$, and $|\mathsf Z'_{p-2j,N}|\le C$ on
$[-1,1]$ by \cref{lem:ent-harmonic}(3), and the term is bounded. There are
finitely many terms, and $|r|^{p-1}\le|r|^{d-1}$ since $p\ge d$. This proves
the bound on $m_x(r)$.
\end{proof}

We now define the truncation. We fix once and for all an even
$C^\infty$ function $\chi:\R\to[0,1]$ with $\chi=1$ on $[-1,1]$ and $\chi=0$
outside $[-2,2]$. For a cutoff scale $\alpha>0$, $u,x\in S_N$, and an
admissible pair $(h,k)$, we set
\begin{equation}\label{eq:ent-cutoff}
 \begin{gathered}
 r=R(x,u),\qquad s=\frac{G_u\cdot x}{N},\\
 \delta_\alpha(u,x)=\chi(r/\alpha)\chi(s/\alpha)\,\delta_{h,k,u}(x),\qquad
 w(u,x)=\beta\delta_\alpha(u,x).
 \end{gathered}
\end{equation}
The function $\delta_\alpha(u,\cdot)$ is smooth on $\R^N$, and it depends on
the field through $H_Q(u)$, $K_Q(u)$, and $G_u$.

\begin{proposition}[Axis-averaged moments of the density ratio]\label{prop:ent-multipliers}
There exist $\alpha_0\in(0,1/8]$, $C_1$, and $N_0$ such that, for every
$\alpha\in(0,\alpha_0]$, every $N\ge N_0$, every field in $\Omega_N$, every
admissible $(h,k)$, and every $x\in S_N$, with $u$ uniform on $S_N$,
\begin{align}
 \E_u\bigl|e^{w(u,x)}-1\bigr|&\le C_1N^{1-d/2},
 \label{eq:ent-eta}\\
 \E_u\bigl[e^{w(u,x)}\bigl(1-e^{-w(u,x)}\bigr)^2\bigr]&\le C_1N^{2-d}.
 \label{eq:ent-epsilon}
\end{align}
Moreover, for all $u\in S_N$,
\begin{equation}\label{eq:ent-cutoff-sup}
 \sup_{x\in S_N}|\delta_\alpha(u,x)|\le C_1N\alpha^d,
 \qquad
 \norm{\nabla\delta_\alpha(u,\cdot)}_\infty\le C_1\sqrt N\alpha^{d-1}.
\end{equation}
The constants $\alpha_0,C_1,N_0$ do not depend on $\alpha$.
\end{proposition}

\begin{proof}
The proof has three steps. We first bound exponential moments on each
latitude, then integrate over the overlap $r$, and finally verify the
deterministic bounds.

\paragraph{Moments on a latitude.}
Fix $\zeta\ge0$; we will take $\zeta=2\beta$. Recall that $|G_u\cdot x|=\sqrt N|V_x(u)|$.
Using \eqref{eq:ent-delta-bounds}, we find
\begin{equation}\label{eq:ent-delta-V}
 |\delta_{h,k,u}(x)|\le C_2N|r|^d+C_2\sqrt N|r|^{d-1}|V_x(u)|.
\end{equation}
Here $V_x$ is the gradient of the sampled field itself. A centered variable
with the bound \eqref{eq:ent-latitude} satisfies
$\Pp(|Z|>v)\le2e^{-v^2/(4C)}$ by Chernoff's inequality. It follows that
$\E_{\rm lat}Z^4\le C$ and
$\E_{\rm lat}e^{t|Z|}\le2e^{Ct^2}$ for $t\ge0$. The Cauchy--Schwarz
inequality then gives $\E_{\rm lat}[Z^2e^{t|Z|}]\le Ce^{2Ct^2}$. We write
$V_x=m_x+Z$, insert the bound on $m_x$ from \eqref{eq:ent-latitude} into
\eqref{eq:ent-delta-V}, and apply these estimates with
$t=\zeta C_2\sqrt N|r|^{d-1}$. We obtain
\begin{equation}\label{eq:ent-lat-moments}
 \E_{\rm lat}e^{\zeta|\delta|}\le C_\zeta e^{\Psi_\zeta(r)},\qquad
 \E_{\rm lat}\bigl[V_x^2e^{\zeta|\delta|}\bigr]
 \le C_\zeta\bigl(1+N|r|^{2d-2}\bigr)e^{\Psi_\zeta(r)},
\end{equation}
where $\delta=\delta_{h,k,u}(x)$ and
\[
 \Psi_\zeta(r)=C_\zeta\bigl(N|r|^d+N|r|^{2d-2}+\sqrt N|r|^{d-1}\bigr).
\]
In $\Psi_\zeta$, the first term comes from the value correction, the second
from the square of the random coefficient $t$ and from the mean $m_x$, and
the third from the bounded part of $m_x$.

\paragraph{Integration over the overlap.}
For $N\ge6$, the density \eqref{eq:ent-axis-density} satisfies
$p_N(r)\le\sqrt Ne^{-Nr^2/4}$. We choose $\alpha_0\le1/8$ so small that
$C_\zeta\bigl((2\alpha_0)^{d-2}+(2\alpha_0)^{2d-4}+(2\alpha_0)^{d-3}\bigr)
\le1/8$. Since $\sqrt N\le N$ and $d\ge4$, we then have
$\Psi_\zeta(r)\le Nr^2/8$ for $|r|\le2\alpha_0$. For every fixed $j\ge0$ and
$\alpha\le\alpha_0$, it follows that
\begin{equation}\label{eq:ent-overlap-moments}
 \int_{|r|\le2\alpha}p_N(r)|r|^je^{\Psi_\zeta(r)}\dd r
 \le\sqrt N\int_\R|r|^je^{-Nr^2/8}\dd r\le C_jN^{-j/2}.
\end{equation}
Note that $\delta_\alpha^2e^{\zeta|\delta_\alpha|}\le
\ind_{\{|r|\le2\alpha\}}\delta^2e^{\zeta|\delta|}$, since the cutoffs take
values in $[0,1]$ and vanish for $|r|>2\alpha$. Squaring \eqref{eq:ent-delta-V}, taking the latitude average
with \eqref{eq:ent-lat-moments}, and integrating with
\eqref{eq:ent-overlap-moments}, we find
\begin{equation}\label{eq:ent-weighted-second}
 \E_u\bigl[\delta_\alpha^2e^{\zeta|\delta_\alpha|}\bigr]
 \le C\bigl(N^2N^{-d}+N\,N^{-(d-1)}+N^2N^{-(2d-2)}\bigr)\le CN^{2-d}.
\end{equation}
The three terms come from $N^2r^{2d}$, from $Nr^{2d-2}$ times the bounded
part of the second moment of $V_x$, and from $Nr^{2d-2}$ times
$Nr^{2d-2}$.

We now take $\zeta=2\beta$. The inequalities $|e^t-1|\le|t|e^{|t|}$ and
$e^t(1-e^{-t})^2=e^t+e^{-t}-2\le t^2e^{|t|}$, together with the
Cauchy--Schwarz inequality, give
\[
 \E_u|e^w-1|\le\beta\bigl(\E_u[\delta_\alpha^2e^{2\beta|\delta_\alpha|}]
 \bigr)^{1/2},\qquad
 \E_u[e^w(1-e^{-w})^2]\le\beta^2\E_u[\delta_\alpha^2e^{2\beta|\delta_\alpha|}].
\]
Together with \eqref{eq:ent-weighted-second}, this proves \eqref{eq:ent-eta} and
\eqref{eq:ent-epsilon}.

\paragraph{Deterministic bounds.}
Wherever the cutoff product or one of its derivatives is nonzero, we have
$|r|\le2\alpha$ and $|s|\le2\alpha$. Since $G_u\cdot x=Ns$,
\eqref{eq:ent-theta-orders} and the bounds on $\Omega_N$ imply that
$|\delta_{h,k,u}|\le CN\alpha^d$ there. In
\eqref{eq:ent-correction-gradient}, the terms with $\theta_1',\theta_2'$
cost $C\sqrt N\alpha^{d-1}$, the term with $\theta_3'$ costs
$C\alpha^{d-2}N|s|N^{-1/2}\le C\sqrt N\alpha^{d-1}$, and the last term costs
$C\sqrt N\alpha^{d-1}$. Further, $|\nabla r|\le N^{-1/2}$ and
$|\nabla s|=|\Pi_xG_u|/N\le C_0N^{-1/2}$, and each differentiated cutoff
contributes at most $C\alpha^{-1}N^{-1/2}\cdot N\alpha^d$. These estimates
prove \eqref{eq:ent-cutoff-sup}.
\end{proof}

\section{An entropy--resolvent comparison}\label{sec:ent-resolvent}

The next lemma does not involve the Gaussian model. It compares the killing
time of a perturbed reversible diffusion, killed at a rate $\kappa$, with the
integrated killing rate against a resolvent of the unperturbed diffusion. The perturbation is
averaged over a frozen label, and only the two averaged moments of
\eqref{eq:ent-eta} and \eqref{eq:ent-epsilon} enter. We state the lemma for
diffusion Dirichlet forms, so that integration by parts with frozen
variables and bounded measurable perturbations has a precise meaning. Its
ingredients are standard Dirichlet form calculus \cite{FukushimaOshimaTakeda2011}; the quantitative comparison below is proved here.

\paragraph{Setting.}
Let $(\mathcal U,\vartheta)$ and $(E,\mu)$ be probability spaces and
$\bar\mu=\vartheta\otimes\mu$. We think of $u\in\mathcal U$ as a frozen
label, that is, a parameter that does not move under the dynamics. Let $\cE_0$ be a symmetric Dirichlet form on $L^2(\bar\mu)$ with domain
$\mathcal V$, and let $A_0\ge0$ be the associated self-adjoint operator, so
that $\cE_0(g,v)=\ip{A_0g}{v}_{\bar\mu}$ for $g\in\operatorname{Dom}(A_0)$
and $v\in\mathcal V$. We assume the following.
\begin{enumerate}
\item[(D1)] (Conservativeness.) $1\in\mathcal V$ and $\cE_0(1,1)=0$.
\item[(D2)] (Carr\'e du champ.) There exists a symmetric bilinear map
$(g,v)\mapsto Dg\cdot Dv$ from $\mathcal V\times\mathcal V$ to
$L^1(\bar\mu)$, with $Dg\cdot Dg\ge0$, such that
$\cE_0(g,v)=\int Dg\cdot Dv\dd\bar\mu$.
\item[(D3)] (Diffusion property.) If $g,v\in\mathcal V\cap L^\infty$, then
$gv\in\mathcal V$ and $D(gv)\cdot D\psi=g\,Dv\cdot D\psi+v\,Dg\cdot D\psi$
for $\psi\in\mathcal V$. If $g\in\mathcal V\cap L^\infty$ and
$\varphi\in C^1(\R)$, then $\varphi(g)\in\mathcal V$ and
$D\varphi(g)\cdot D\psi=\varphi'(g)\,Dg\cdot D\psi$.
\item[(D4)] (Frozen label.) The resolvents $(\lambda+A_0)^{-1}$,
$\lambda>0$, and the map $(g,v)\mapsto Dg\cdot Dv$ send functions that do
not depend on $u$ to functions that do not depend on $u$.
\end{enumerate}
By (D2), $|Dg\cdot Dv|\le(Dg\cdot Dg)^{1/2}(Dv\cdot Dv)^{1/2}$ pointwise. We
write $|Dg|^2=Dg\cdot Dg$, and identify a function on $E$ with a function on
$\mathcal U\times E$ that does not depend on $u$. In our application, $E$ is
a product of a parameter space and $S_N$, the label is the axis, and $Dg\cdot
Dv=\nabla g\cdot\nabla v$ with the gradient acting in the variable of $S_N$
only.

Let $w$ be a bounded measurable function on $\mathcal U\times E$. We define
\[
 M=\int e^w\dd\bar\mu,\qquad \mu_w=M^{-1}e^w\bar\mu,\qquad
 \cE_w(g,v)=\int Dg\cdot Dv\dd\mu_w,\quad g,v\in\mathcal V .
\]
Since $w$ is bounded, the norms of $L^2(\mu_w)$ and $L^2(\bar\mu)$ are
equivalent, and $\cE_w$ is a closed form on $L^2(\mu_w)$ with domain
$\mathcal V$; let $A_w\ge0$ be its operator. For a smooth $w$ on a manifold,
$-A_w=-A_0+Dw\cdot D$. We assume that there exists a Markov process
$(u,X_t)_{t\ge0}$ on $\mathcal U\times E$, whose label coordinate $u$ does not
move, and whose transition semigroup acts on bounded measurable functions by
$e^{-tA_w}$, $\mu_w$-almost everywhere.

\begin{lemma}[Entropy--resolvent comparison]\label{lem:ent-resolvent}
In the setting above, suppose that for $\mu$-almost every $x\in E$,
\begin{equation}\label{eq:ent-abstract-moments}
 \int\bigl|e^{w(u,x)}-1\bigr|\dd\vartheta(u)\le\eta\le\frac12,
 \qquad
 \int e^{w(u,x)}\bigl(1-e^{-w(u,x)}\bigr)^2\dd\vartheta(u)\le\epsilon .
\end{equation}
Let $f_0\ge0$ be a function on $E$ with $\int f_0\dd\mu=1$ and
$\mathcal H_0=\int(1+f_0)\log(1+f_0)\dd\mu<\infty$. Let $\kappa\ge0$ be a
bounded measurable function on $\mathcal U\times E$. Start the process
from the law $\vartheta\otimes(f_0\mu)$, let $\mathbf e$ be a unit
exponential variable independent of it, and let
$\tau_\kappa=\inf\{t\ge0:\int_0^t\kappa(u,X_v)\dd v\ge\mathbf e\}$, with
$\inf\emptyset=\infty$, be the killing time at rate $\kappa$. For $\lambda>0$ define
\begin{equation}\label{eq:ent-resolvents}
 f=\lambda(\lambda+A_0)^{-1}f_0,\qquad F=1+f,\qquad
 \mathcal R_\lambda=\lambda^{-1}\int\kappa F\dd\mu_w .
\end{equation}
Then
\begin{equation}\label{eq:ent-resolvent-bound}
 \E e^{-\lambda\tau_\kappa}\le4\eta+5\sqrt\epsilon+\mathcal R_\lambda
 +2\sqrt{\epsilon\mathcal H_0\bigl[\mathcal R_\lambda
   +\epsilon(1+\mathcal H_0)\bigr]},
\end{equation}
with $e^{-\lambda\cdot\infty}=0$. In particular, for every $T>0$,
\begin{equation}\label{eq:ent-resolvent-horizon}
 \Pp(\tau_\kappa\le T)\le14\bigl(\eta+\sqrt\epsilon
 +\epsilon(1+\mathcal H_0)+\mathcal R_{1/T}\bigr).
\end{equation}
\end{lemma}

The first resolvent in \eqref{eq:ent-resolvents} acts on the unperturbed
form; for unbounded $f_0$ it is extended to $L^1(\mu)$, where it is a
contraction. No constant in \eqref{eq:ent-resolvent-bound} depends on
$\norm{w}_\infty$, on $\lambda$, or on a spectral gap, and the form need not
be irreducible. The boundedness of $w$ is used only to define $\cE_w$ and
$A_w$.

\begin{proof}
We first treat bounded $f_0$. An approximation removes this restriction at
the end. The proof has four steps: an entropy bound for the resolvent of the
unperturbed generator, the probabilistic meaning of the resolvent of the
killed diffusion, a weighted energy
estimate, and a comparison of initial laws.

\paragraph{Two consequences of the moment bounds.}
By the first bound in \eqref{eq:ent-abstract-moments}, we have
$1-\eta\le M\le1+\eta$. Let $k=1-e^{-w}$. For every nonnegative measurable
$\Gamma$ on $E$, Tonelli's theorem and \eqref{eq:ent-abstract-moments} give
\begin{equation}\label{eq:ent-k-square}
 \int k^2\Gamma\dd\mu_w\le2\epsilon\int\Gamma\dd\mu,\qquad
 \int\Gamma\dd\mu_w\le3\int\Gamma\dd\mu .
\end{equation}
Indeed, $e^wk^2=e^w(1-e^{-w})^2$, $M\ge1/2$, and $(1+\eta)/(1-\eta)\le3$.
All the small constants below come from \eqref{eq:ent-k-square}, not from a
pointwise bound on $e^w$.

\paragraph{The resolvent of the unperturbed generator.}
For bounded $f_0$, the representation
$f=\lambda\int_0^\infty e^{-\lambda t}e^{-tA_0}f_0\dd t$ and (D1) show that
$f\ge0$, $\int f\dd\mu=1$, and $f\in\operatorname{Dom}(A_0)\cap L^\infty$.
By (D4), $f$ does not depend on $u$. With $F_0=1+f_0$, we have
$A_0F=A_0f=\lambda(F_0-F)$. The function $F\ge1$ is bounded, and (D3) shows that
$\log F\in\mathcal V$ with $D\log F\cdot D\psi=F^{-1}DF\cdot D\psi$. Let
$\mathcal J=\int|DF|^2F^{-1}\dd\mu$. Testing with $\log F$, we find
\begin{equation}\label{eq:ent-fisher}
 \begin{split}
 \mathcal J&=\ip{A_0F}{\log F}_\mu=\lambda\ip{F_0-F}{\log F}_\mu\\
 &\le\lambda\Bigl[\int F_0\log F_0\dd\mu-\int F\log F\dd\mu\Bigr]
 \le\lambda\mathcal H_0 .
 \end{split}
\end{equation}
In the first inequality, we used the tangent line inequality
$(a-b)\log b\le a\log a-b\log b-(a-b)$ for $a,b>0$, and that $F_0$ and $F$
both have integral two. In the second, we used $F\log F\ge0$.

\paragraph{The resolvent of the killed diffusion.}
The form $\cE_w(g,v)+\int\kappa gv\dd\mu_w$ is closed on $\mathcal V$; let
\[
 h=(\lambda+A_w+\kappa)^{-1}\kappa\in\operatorname{Dom}(A_w)\subset
 \mathcal V .
\]
The Duhamel formula expresses $e^{-t(A_w+\kappa)}$
as the series
\[
 e^{-t(A_w+\kappa)}=\sum_{n\ge0}(-1)^n\int_{0<v_1<\dots<v_n<t}
 e^{-v_1A_w}\kappa e^{-(v_2-v_1)A_w}\kappa\cdots\kappa
 e^{-(t-v_n)A_w}\dd v,
\]
whose $n$th term has norm at most $\norm{\kappa}_\infty^nt^n/n!$ on bounded
functions.
Expanding the exponential in
$\E_{u,x}\bigl[e^{-\int_0^t\kappa(u,X_v)\dd v}g(X_t)\bigr]$ and using the
Markov property, we identify the two series term by term. The result is
the Feynman--Kac formula
$e^{-t(A_w+\kappa)}g(u,x)=\E_{u,x}\bigl[e^{-\int_0^t\kappa(u,X_v)\dd v}
g(X_t)\bigr]$ for bounded $g$. Conditionally on the path, $\tau_\kappa$ has
density $\kappa(u,X_t)e^{-\int_0^t\kappa(u,X_v)\dd v}$ on $[0,\infty)$.
Integrating the Feynman--Kac formula for $g=\kappa$ against
$e^{-\lambda t}\dd t$, we obtain
\begin{equation}\label{eq:ent-hazard-transform}
 h(u,x)=\E_{u,x}e^{-\lambda\tau_\kappa}.
\end{equation}
Note that $0\le h\le1$. This is why the lemma controls the killing
time itself.

\paragraph{The commutator.}
Let $g\in\operatorname{Dom}(A_0)$ not depend on $u$, and let $v\in\mathcal
V$ be bounded. Then
\begin{equation}\label{eq:ent-commutator}
 \cE_w(g,v)=\ip{(1-k)A_0g}{v}_{\mu_w}+\int k\,Dg\cdot Dv\dd\mu_w .
\end{equation}
Indeed, since $(1-k)e^w=1$, the first term on the right equals
$M^{-1}\int A_0g\,v\dd\bar\mu=M^{-1}\int Dg\cdot Dv\dd\bar\mu$, and adding the
second term multiplies the integrand by $1+(e^w-1)=e^w$. The identity
\eqref{eq:ent-commutator} is a weak form of $A_wg=(1-k)A_0g+D^*_w(kDg)$,
where $D^*_w$ is the adjoint in $L^2(\mu_w)$; no derivative of $w$ or $k$
appears in it.

\paragraph{The weighted energy.}
Let $B_h=\int F|Dh|^2\dd\mu_w$ and $A_h=\int Fh^2\dd\mu_w$, which are finite.
By (D3), $Fh\in\mathcal V$ and $D(Fh)\cdot Dh=F|Dh|^2+h\,DF\cdot Dh$.
Testing the equation $(\lambda+A_w+\kappa)h=\kappa$ against $Fh$, and using
$h\le1$, we find
\begin{equation}\label{eq:ent-weighted-energy}
 B_h+\int h\,DF\cdot Dh\dd\mu_w+\int(\lambda+\kappa)Fh^2\dd\mu_w
 =\int\kappa Fh\dd\mu_w\le\int\kappa F\dd\mu_w=\lambda\mathcal R_\lambda .
\end{equation}
The cross term is the delicate one, since $|DF|$ is not small. Using (D3), we write
$h\,DF\cdot Dh=\frac12DF\cdot D(h^2)$. Applying \eqref{eq:ent-commutator}
with $g=F$ and $v=h^2$, and using $A_0F=\lambda(f_0-f)$ and
$(1-k)e^w=1$, we obtain
\begin{equation}\label{eq:ent-cross}
 \int h\,DF\cdot Dh\dd\mu_w
 =\frac{\lambda}{2M}\int f_0h^2\dd\bar\mu
 -\frac\lambda2\int fh^2\dd\mu_w+\frac\lambda2\int kfh^2\dd\mu_w
 +\int kh\,DF\cdot Dh\dd\mu_w .
\end{equation}
The first term on the right is nonnegative. Since $F-f/2\ge F/2$, the second
term combined with $\lambda A_h$ in \eqref{eq:ent-weighted-energy} leaves at
least $\lambda A_h/2$. We apply \eqref{eq:ent-k-square} to the third and fourth
terms, with $\Gamma=f$ and $\Gamma=|DF|^2/F$, which do not depend on $u$ by
(D4). Using also $h\le1$, $\int f\dd\mu=1$, and the Cauchy--Schwarz
inequality, we obtain
\begin{equation}\label{eq:ent-energy-errors}
 \Bigl|\int kfh^2\dd\mu_w\Bigr|\le\sqrt{2\epsilon A_h},\qquad
 \Bigl|\int kh\,DF\cdot Dh\dd\mu_w\Bigr|\le\sqrt{2\epsilon\mathcal JB_h}.
\end{equation}
For the first bound, we used $\int fh^4\dd\mu_w\le A_h$. Inserting
\eqref{eq:ent-cross} and \eqref{eq:ent-energy-errors} into
\eqref{eq:ent-weighted-energy}, we have
\[
 B_h+\frac\lambda2A_h\le\lambda\mathcal R_\lambda
 +\frac\lambda2\sqrt{2\epsilon A_h}+\sqrt{2\epsilon\mathcal JB_h}.
\]
Using Young's inequality in the forms $\frac\lambda2\sqrt{2\epsilon A_h}
\le\frac\lambda4A_h+\frac\lambda2\epsilon$ and
$\sqrt{2\epsilon\mathcal JB_h}\le\frac12B_h+\epsilon\mathcal J$, and then
\eqref{eq:ent-fisher}, we find
\begin{equation}\label{eq:ent-energy}
 \frac{B_h}{\lambda}\le2\mathcal R_\lambda+\epsilon+2\epsilon\mathcal H_0 .
\end{equation}
The term $\lambda A_h$ must be kept here, since it absorbs the negative
second term of \eqref{eq:ent-cross}. Without it, $B_h/\lambda$ would only be
bounded by a constant, and the last term of \eqref{eq:ent-resolvent-bound}
below would be of order $\sqrt{\epsilon\mathcal H_0}$, which is
$N^{(3-d)/2}$ in our application.

\paragraph{Comparison of initial laws.}
The actual initial law $\vartheta\otimes(f_0\mu)$ has density
$q_0=Me^{-w}f_0$ with respect to $\mu_w$. Because the first moment bound in
\eqref{eq:ent-abstract-moments} holds for each $x$, we have
\begin{equation}\label{eq:ent-initial}
 \norm{q_0-f_0}_{L^1(\mu_w)}=M^{-1}\int f_0|M-e^w|\dd\bar\mu
 \le2\int f_0\Bigl(|M-1|+\int|1-e^w|\dd\vartheta\Bigr)\dd\mu\le4\eta .
\end{equation}
By \eqref{eq:ent-hazard-transform}, $\E e^{-\lambda\tau_\kappa}
=\ip{q_0}{h}_{\mu_w}\le\ip{f_0}{h}_{\mu_w}+4\eta$. To bound the last pairing,
we test the equation for $h$ against $f$. Using \eqref{eq:ent-commutator} with $g=f$ and $v=h$, we have $\cE_w(f,h)=\lambda\ip{f_0-f}{h}_{\mu_w}
-\ip{kA_0f}{h}_{\mu_w}+\int k\,Df\cdot Dh\dd\mu_w$. Inserting this
into the equation
$\lambda\ip{h}{f}_{\mu_w}+\cE_w(h,f)+\ip{\kappa h}{f}_{\mu_w}
=\ip{\kappa}{f}_{\mu_w}$, we obtain
\begin{equation}\label{eq:ent-source}
 \ip{f_0}{h}_{\mu_w}=\lambda^{-1}\ip{f}{\kappa(1-h)}_{\mu_w}
 +\ip{k(f_0-f)}{h}_{\mu_w}-\lambda^{-1}\int k\,Df\cdot Dh\dd\mu_w .
\end{equation}
The first term lies in $[0,\mathcal R_\lambda]$. For the second, by
the Cauchy--Schwarz inequality and \eqref{eq:ent-k-square} with $\Gamma=f_0$
and $\Gamma=f$, we have
$\int|k|(f_0+f)\dd\mu_w\le2\sqrt{6\epsilon}\le5\sqrt\epsilon$.
For the third, note that $Df=DF$. The argument for
\eqref{eq:ent-energy-errors}, followed by \eqref{eq:ent-fisher} and
\eqref{eq:ent-energy}, yields
\[
 \lambda^{-1}\Bigl|\int k\,Df\cdot Dh\dd\mu_w\Bigr|
 \le\lambda^{-1}\sqrt{2\epsilon\mathcal JB_h}
 \le\sqrt{2\epsilon\mathcal H_0\,B_h/\lambda}
 \le2\sqrt{\epsilon\mathcal H_0\bigl[\mathcal R_\lambda
   +\epsilon(1+\mathcal H_0)\bigr]}.
\]
Combining these bounds proves \eqref{eq:ent-resolvent-bound} for bounded
$f_0$. All the integrations by parts above are identities between forms, for
the functions $F$, $h$, $Fh$, $h^2$, and $\log F$, which lie in $\mathcal V$
by (D3). We use no regularity of these functions or of $w$ beyond
membership in $\mathcal V$.

\paragraph{Unbounded initial densities.}
For general $f_0$, let $f_0^{(n)}=(f_0\wedge n)/\int(f_0\wedge n)\dd\mu$.
Then $f_0^{(n)}\to f_0$ in $L^1(\mu)$. For large $n$, we have
$f_0^{(n)}\le2f_0$. Since $(1+2a)\log(1+2a)\le2(1+a)\log(1+a)+2a$ for
$a\ge0$, dominated convergence gives $\mathcal H_0^{(n)}\to\mathcal H_0$. The
resolvent $\lambda(\lambda+A_0)^{-1}$ is a contraction on $L^1(\mu)$, which implies that $f^{(n)}\to f$
in $L^1(\mu)$. Since $\kappa$ and $e^w$ are bounded, it follows that
$\mathcal R_\lambda^{(n)}\to\mathcal R_\lambda$. The function $h$ does not depend on $f_0$ and is bounded,
so the left side of \eqref{eq:ent-resolvent-bound} converges as well. This
proves \eqref{eq:ent-resolvent-bound} in general.

\paragraph{The time horizon.}
Finally, $\ind_{\{\tau_\kappa\le T\}}\le e\cdot e^{-\tau_\kappa/T}$. Using
\eqref{eq:ent-resolvent-bound} with $\lambda=1/T$, together with
$2\sqrt{\epsilon\mathcal H_0\mathcal R}\le\mathcal R+\epsilon\mathcal H_0$ and
$2\epsilon\sqrt{\mathcal H_0(1+\mathcal H_0)}\le2\epsilon(1+\mathcal H_0)$, we
find that $\Pp(\tau_\kappa\le T)\le e\bigl(4\eta+5\sqrt\epsilon
+2\mathcal R_{1/T}+3\epsilon(1+\mathcal H_0)\bigr)$. This implies
\eqref{eq:ent-resolvent-horizon}.
\end{proof}

\section{Killing-time bound for the truncated process}\label{sec:ent-truncated}

We now apply \cref{lem:ent-resolvent} to the diffusion whose Hamiltonian is
$H_N+\delta_\alpha(u,\cdot)$, with $u$ uniform and independent of the field.
Recall the constants $\alpha_0$ and $C_1$ of \cref{prop:ent-multipliers}.

For a smooth function $\widetilde H$ on $\R^N$, the diffusion with generator
$L_{\beta,\widetilde H}$ is the solution of the stochastic differential
equation
\begin{equation}\label{eq:ent-sde}
 \dd X_t=\sqrt2\,\Pi_{X_t}\dd B_t
 +\Bigl[\beta \Pi_{X_t}\nabla_E\widetilde H(X_t)
 -\frac{N-1}{N}X_t\Bigr]\dd t
\end{equation}
in $\R^N$, where $B$ is a standard Brownian motion in $\R^N$; compare
\cref{lem:pre-dynamics} and \eqref{eq:red-sde}. The coefficients are smooth and in particular
locally Lipschitz. Let $\rho_t=|X_t|^2-N$. Note that $x\cdot \Pi_xv=-\rho\,x\cdot v/N$
and $\operatorname{tr}(\Pi_x\Pi_x)=N-1+\rho^2/N^2$ for $\rho=|x|^2-N$.
Then It\^o's formula gives $\dd\rho_t=\rho_t\,(a_t\dd t+b_t\cdot\dd B_t)$
with locally bounded processes $a_t$ and $b_t$. This implies that $\rho_t=0$
for all $t$ when $\rho_0=0$, so that the solution started on $S_N$ is
global and pathwise unique. When the coefficients of
$\widetilde H$ range over a finite-dimensional parameter space, localization
on bounded sets of parameters and successive approximations construct the
solution as a jointly measurable function of the parameters, the initial
point, and the path of $B$.

\begin{proposition}[The truncated process avoids the caps]
\label{prop:ent-truncated}
There exists $\alpha_1\in(0,\alpha_0]$ such that, for every
$\alpha\in(0,\alpha_1]$, there exist $\gamma_0,c,C>0$ and $N_0$ with the
following property. Let $\mathsf Q_\alpha$ be the joint law of: the
coefficient vector $J$ of $H_N$, conditioned on $\Omega_N$; an independent
uniform axis $u\in S_N$; an independent initial point $X_0$ with law
$\nu_N$; and the solution $X$ of \eqref{eq:ent-sde} with
$\widetilde H=H_J+\delta_\alpha(u,\cdot)$, driven by an independent Brownian
motion. Then, for every $N\ge N_0$ and every admissible $(h,k)$,
\begin{equation}\label{eq:ent-truncated}
 \mathsf Q_\alpha\Bigl(\sup_{0\le t\le e^{\gamma_0N}}|R(X_t,u)|
 \ge\frac\alpha2\Bigr)\le C\bigl(N^{1-d/2}+N^{3-d}\bigr)+e^{-cN}.
\end{equation}
The same bound holds with $\gamma_0$ replaced by any $\gamma\in(0,\gamma_0]$.
\end{proposition}

\begin{proof}
The proof has three steps: we verify the hypotheses of
\cref{lem:ent-resolvent}, bound the integrated killing rate, and convert the
killing at a smooth rate into a hit of the cap $C_{\alpha/2}(u)$.

\paragraph{The setting of the comparison lemma.}
Let $\Theta$ be the space of coefficient vectors, $\Pp_\Omega(\dd J)
=\Pp(\dd J\mid\Omega_N)$, $E=\Theta\times S_N$, and
\[
 \mu(\dd J,\dd x)=\Pp_\Omega(\dd J)\,\pi^{H_J}(\dd x),\qquad
 f_0(J,x)=Z_{H_J}e^{-\beta H_J(x)},
\]
where $Z_{H_J}=\int e^{\beta H_J}\dd\nu_N$. The label space is
$(\mathcal U,\vartheta)=(S_N,\nu_N)$, the space of axes. Then $f_0\mu$ is
exactly $\Pp_\Omega\otimes\nu_N$, the law of the field and the initial
point. On $\Omega_N$, we have $|H_J|\le C_0N$, which implies
$f_0\le e^{2\beta C_0N}$ and
\begin{equation}\label{eq:ent-start-entropy}
 \mathcal H_0\le\log\bigl(1+e^{2\beta C_0N}\bigr)\int(1+f_0)\dd\mu
 \le C(1+N).
\end{equation}
The base form is $\cE_0(g,v)=\int\nabla_xg\cdot\nabla_xv\dd\bar\mu$, where
$\bar\mu=\nu_N\otimes\mu$ and $\nabla_x$ is the gradient in the variable
of $S_N$, with domain
$\mathcal V=L^2\bigl(\nu_N\otimes\Pp_\Omega;H^1(S_N,\nu_N)\bigr)$. Since
$e^{-2\beta C_0N}\le\dd\mu/\dd(\Pp_\Omega\otimes\nu_N)\le e^{2\beta C_0N}$,
the form norm $\bigl(\norm{g}_{L^2(\bar\mu)}^2+\cE_0(g,g)\bigr)^{1/2}$ on
$\mathcal V$ is equivalent to the norm of this Hilbert space. Then $\cE_0$ is closed. Constants lie in $\mathcal V$ with zero
energy, the product and chain rules hold fiberwise, and the operations in
(D4) act fiberwise in $(u,J)$. The equivalence constants grow with $N$, but
they are never used as small constants. For fixed $(u,J)$, the operator
$A_0$ acts as $-L_{\beta,H_J}$, which is symmetric for $\pi^{H_J}$.

We take $w=\beta\delta_\alpha(u,x)$, a bounded function of $(u,J,x)$, smooth
in $x$. The conditional law of $\mu_w$ given $(u,J)$ is the Gibbs measure
of $H_J+\delta_\alpha(u,\cdot)$. For fixed $(u,J)$, the form $\cE_w$ is
therefore the Dirichlet form of $L_{\beta,H_J+\delta_\alpha(u,\cdot)}$, and $A_w$ acts as
$-L_{\beta,H_J+\delta_\alpha(u,\cdot)}$. The solutions of \eqref{eq:ent-sde},
constructed for each fixed $(u,J)$ and measurable in these parameters,
provide the Markov process required by the lemma, with frozen label $u$ and
frozen parameter $J$. Its initial law $\vartheta\otimes(f_0\mu)$ is the
initial law of $\mathsf Q_\alpha$. No differentiation in $J$ is needed,
since the gradient acts only in $x$.

For $\alpha\le\alpha_0$, \cref{prop:ent-multipliers} verifies
\eqref{eq:ent-abstract-moments} for every $(J,x)$ with $J\in\Omega_N$, with
\begin{equation}\label{eq:ent-applied-moments}
 \eta=C_1N^{1-d/2},\qquad\epsilon=C_1N^{2-d},
\end{equation}
and $\eta\le1/2$ for large $N$.

\paragraph{The integrated killing rate.}
Let $\varphi:\R\to[0,1]$ be smooth and even, with $\varphi(r)=0$ for
$|r|\le\alpha/4$ and $\varphi(r)=1$ for $|r|\ge\alpha/3$, and let
$\kappa(u,J,x)=\varphi(R(x,u))$. Since $f$ does not depend on $u$ and
$\int F\dd\mu=2$, we integrate over the label first, bound $e^w$ by
\eqref{eq:ent-cutoff-sup}, and use $M\ge1/2$. For fixed $x$, the set of
axes $u$ with $\kappa(u,J,x)\neq0$ is contained in $C_{\alpha/4}(x)$. By
symmetry, $\nu_N(C_b(x))$ does not depend on $x\in S_N$; we denote this
common value by $\nu_N(C_b)$. We obtain
\begin{equation}\label{eq:ent-hazard-mass}
 \mathcal R_\lambda\le4\lambda^{-1}e^{C_1\beta N\alpha^d}\,
 \nu_N\bigl(C_{\alpha/4}\bigr).
\end{equation}
This step uses that \eqref{eq:ent-cutoff-sup} holds for every $u$ and $x$,
not only on average. By \eqref{eq:ent-axis-density}, for every fixed
$b\in(0,1)$ and large $N$,
\begin{equation}\label{eq:ent-cap}
 \nu_N(C_b)=2\int_b^1p_N(r)\dd r\le2\sqrt N\,e^{-(N-3)I(b)}
 \le e^{-NI(b)/2}.
\end{equation}
Since $I(b)\ge b^2/2$ and $d>2$, there exists $\alpha_1\in(0,\alpha_0]$
such that $C_1\beta\alpha^d\le I(\alpha/4)/8$ for all $\alpha\le\alpha_1$.
We fix $\alpha\le\alpha_1$ and set $\gamma_0=I(\alpha/4)/8$. For
$\lambda=(e^{\gamma_0N}+1)^{-1}$, \eqref{eq:ent-hazard-mass} and
\eqref{eq:ent-cap} imply that $\mathcal R_\lambda\le8e^{-NI(\alpha/4)/4}$
for large $N$. Combining \eqref{eq:ent-resolvent-horizon}, \eqref{eq:ent-start-entropy},
and \eqref{eq:ent-applied-moments}, we obtain
\begin{equation}\label{eq:ent-first-firing}
 \mathsf Q_\alpha\bigl(\tau_\kappa\le e^{\gamma_0N}+1\bigr)
 \le C\bigl(N^{1-d/2}+N^{3-d}\bigr)+e^{-cN},
\end{equation}
where $\tau_\kappa$ is the killing time at rate $\kappa$, defined with an
independent unit exponential variable as in \cref{lem:ent-resolvent}. In
\eqref{eq:ent-first-firing} we used $N^{2-d}\le N^{1-d/2}$.

\paragraph{From the killing time to the cap.}
Let $r_t=R(X_t,u)$. By \eqref{eq:ent-sde}, $r_t$ has drift
$-\frac{N-1}{N}r_t+\beta N^{-1}u\cdot\nabla\widetilde H(X_t)$ and
martingale part with quadratic variation rate $2|\Pi_{X_t}u|^2/N^2
=2(1-r_t^2)/N\le2/N$. On $\Omega_N$, we have $\norm{\nabla\widetilde H}_\infty\le C\sqrt N$ by
\eqref{eq:ent-regularity} and \eqref{eq:ent-cutoff-sup}, and the drift is
bounded by a constant $\bar d_r$. Let
$\tau_{\alpha/2}=\inf\{t\ge0:|r_t|\ge\alpha/2\}$, let
$\Delta=\alpha/2-\alpha/3=\alpha/6$, and let $t_0=\Delta/(4(\bar d_r+1))<1$. We
claim that, for every $T>0$,
\begin{equation}\label{eq:ent-exposure}
 \mathsf Q_\alpha\Bigl(\{\tau_{\alpha/2}\le T\}\cap
 \Bigl\{\inf_{\tau_{\alpha/2}\le t\le\tau_{\alpha/2}+t_0}|r_t|<\frac\alpha3
 \Bigr\}\Bigr)\le2\exp\Bigl(-\frac{N\Delta^2}{16t_0}\Bigr).
\end{equation}
Indeed, after $\tau_{\alpha/2}$, the drift moves $r_t$ by at most
$\bar d_rt_0\le\Delta/4$ during a time $t_0$. A decrease of $|r_t|$ from at least
$\alpha/2$ to below $\alpha/3$ therefore requires an increment of the
martingale part of absolute value at least $\Delta/2$, and its quadratic
variation during this interval is at most $2t_0/N$. A continuous martingale
whose quadratic variation is at most $v$ has an absolute increment at least
$b$ with probability at most $2e^{-b^2/(2v)}$, by the exponential
supermartingale. Applying this after the stopping time $\tau_{\alpha/2}$,
conditionally on the frozen variables, and integrating, we conclude that
\eqref{eq:ent-exposure} holds. This also covers the case $\tau_{\alpha/2}=0$.

On the complement of the event in \eqref{eq:ent-exposure}, and on
$\{\tau_{\alpha/2}\le T\}$, the killing rate equals one on the whole interval
$[\tau_{\alpha/2},\tau_{\alpha/2}+t_0]$. Conditionally on the path, the
integrated killing rate is at least $t_0$ by time $T+t_0$, and the killing occurs
by that time with probability at least $1-e^{-t_0}$. Integrating, and using \eqref{eq:ent-exposure}, we find
\[
 \mathsf Q_\alpha(\tau_{\alpha/2}\le T)\le\bigl(1-e^{-t_0}\bigr)^{-1}
 \mathsf Q_\alpha(\tau_\kappa\le T+t_0)+2e^{-N\Delta^2/(16t_0)}.
\]
We take $T=e^{\gamma_0N}$ and use \eqref{eq:ent-first-firing} with $t_0<1$.
After adjusting $C$, $c$, and $N_0$, this proves \eqref{eq:ent-truncated}.
The left side of \eqref{eq:ent-truncated} decreases when $\gamma_0$ is
decreased.
\end{proof}

\section{The cutoff in the gradient variable}\label{sec:ent-boundary}

The truncated and true processes coincide until $|r|$ or $|s|$ reaches the
cutoff scale. In this section, we show that the true process is unlikely to
reach the cutoff boundary $|s|=\alpha$. The argument uses the independence of $G_u$
from the conditioned field, under the unrestricted Gaussian law.

\begin{lemma}[The gradient cutoff is not reached]\label{lem:ent-boundary}
Let $(h,k)$ be admissible. Let $u$ be uniform on $S_N$ and independent of
$H_N$, let $X_0$ have law $\nu_N$, and let $X$ solve \eqref{eq:ent-sde} with
$\widetilde H=H^{h,k,u}$, where $X_0$ and the Brownian motion are
independent of $(H_N,u)$. Let $s_t=G_u\cdot X_t/N$. For every $b>0$, there
exist $\gamma_b,c_b>0$ and $N_0$ such that, for $N\ge N_0$,
\begin{equation}\label{eq:ent-boundary}
 \Pp\Bigl(\sup_{0\le t\le e^{\gamma_bN}}|s_t|\ge b\Bigr)\le e^{-c_bN}.
\end{equation}
The same bound, with a smaller $c_b$, holds when the coefficient vector of
$H_N$ is conditioned on $\Omega_N$. The constants do not depend on $(h,k)$.
\end{lemma}

\begin{proof}
We first work under the unrestricted law. Conditionally on $u$, the
coefficient vector of $H_N-D_u$ is independent of $G_u$, by
\cref{prop:ent-regression}, and $G_u$ is a centered Gaussian vector in
$u^\perp$ with covariance $a_QI_{u^\perp}$, by \eqref{eq:ent-covariances}.
The coefficients of $H^{h,k,u}=(H_N-D_u)+N(\theta_1h+\theta_2k)$ are a
continuous function of those of $H_N-D_u$ and of $u$. By the measurable
construction of the solutions of \eqref{eq:ent-sde}, the whole trajectory
$(X_t)_{t\ge0}$ is a measurable function of these coefficients, of $u$, of
$X_0$, and of the Brownian motion. Conditionally on $u$, it is therefore
independent of $G_u$. At each fixed time $t$, conditionally on $u$ and on
the trajectory, $s_t$ is centered Gaussian with variance
$a_Q|\Pi_uX_t|^2/N^2\le a_Q/N$, and
\begin{equation}\label{eq:ent-grid-tail}
 \Pp\bigl(|s_t|\ge b/2\bigr)\le2e^{-Nb^2/(8a_Q)} .
\end{equation}

We choose $0<\gamma_b<b^2/(32a_Q)$ and the grid of times $j/N$,
$0\le j\le\lceil Ne^{\gamma_bN}\rceil$. A union bound in
\eqref{eq:ent-grid-tail} shows that some grid value satisfies
$|s_{j/N}|\ge b/2$ with probability at most $e^{-cN}$. We next control the
intervals between grid points; here we use the regularity event once. On
$\Omega_N$, we have $|G_u|\le C_0\sqrt N$ and $\norm{\nabla H^{h,k,u}}_\infty
\le C\sqrt N$, by \eqref{eq:ent-regularity} and
\eqref{eq:ent-delta-bounds}. By \eqref{eq:ent-sde}, $s_t$ has drift
$-\frac{N-1}{N}s_t+\beta N^{-1}G_u\cdot\nabla H^{h,k,u}(X_t)$, which is
bounded by a constant $\bar d_s$ since $|s_t|\le C_0$, and its martingale part
has quadratic variation rate $2|\Pi_{X_t}G_u|^2/N^2\le\bar q_s/N$. On an
interval of length $1/N$, the drift moves $s_t$ by at most $b/4$ for large
$N$. If $|s|<b/2$ at the left endpoint and $|s|$ reaches $b$ in the interval,
the martingale increment has absolute value at least $b/4$, while its
quadratic variation is at most $\bar q_s/N^2$. The martingale inequality used
in the proof of \cref{prop:ent-truncated} shows that this has probability at most
$2e^{-b^2N^2/(32\bar q_s)}$. This bound holds conditionally on every field in
$\Omega_N$ and every axis, because the Brownian motion is independent of
them. A union over $O(Ne^{\gamma_bN})$ intervals costs at most $e^{-cN^2}$.
Adding $\Pp(\Omega_N^c)\le e^{-N}$ once, we conclude that
\eqref{eq:ent-boundary} holds.

Under the law conditioned on $\Omega_N$, the probability of the event in
\eqref{eq:ent-boundary} is at most its unrestricted probability divided by
$\Pp(\Omega_N)\ge1/2$. This proves the second assertion. The independence of
$G_u$ from the trajectory was used only under the unrestricted law.
\end{proof}

\section{Proof of the main estimate}\label{sec:ent-proof}

\begin{proof}[Proof of \cref{thm:ent-main}]
We fix the parameters, couple the true and truncated processes, remove the
regularity restriction, and finally pass from a uniform axis to a
deterministic one.

\paragraph{Choice of parameters.}
We fix $\xi,Q,\beta,\bar h,\bar k$, and take $C_0$ and $\Omega_N$ from
\cref{lem:ent-regularity}, and $\alpha_1$ from \cref{prop:ent-truncated}. We
set $a_*=\alpha_1/2$. Given $a\in(0,a_*]$, we set $\alpha=2a$, so that the
cap $C_{\alpha/2}(u)$ in \eqref{eq:ent-truncated} is $C_a(u)$. We then choose
\[
 \gamma=\min\{\gamma_0,\gamma_\alpha\},
\]
where $\gamma_0$ is given by \cref{prop:ent-truncated} and $\gamma_\alpha$
by \cref{lem:ent-boundary} with $b=\alpha$. All these choices precede
$N\to\infty$, and none depends on $u$ or $(h,k)$; see
\cref{tab:ent-parameters}.

\paragraph{The coupling.}
Let the coefficient vector $J$ of $H_N$ have its law conditioned on
$\Omega_N$, and sample an independent uniform axis $u$, an independent
initial point $X_0$ with law $\nu_N$, and an independent Brownian motion
$B$. On this probability space, let $X^{\rm tr}$ solve \eqref{eq:ent-sde}
with $\widetilde H=H^{h,k,u}=H_N+\delta_{h,k,u}$, and let $X^{\rm cut}$ solve
it with $\widetilde H=H_N+\delta_\alpha(u,\cdot)$, both from $X_0$ and driven
by $B$. Let $r_t$ and $s_t$ be defined by \eqref{eq:ent-cutoff} along each
path. On the open set $\{x:|R(x,u)|<\alpha,\ |G_u\cdot x|/N<\alpha\}$, the
two cutoffs equal one, and the two Hamiltonians and their gradients agree
there. By pathwise uniqueness, the two solutions coincide up to the first
time that either of them leaves this set.

Let $T=e^{\gamma N}$, let $\tau_a^{\rm tr}$ and $\tau_a^{\rm cut}$ be the
hitting times of $C_a(u)$ by the two processes, and let
$\varsigma^{\rm tr}=\inf\{t\ge0:|s_t^{\rm tr}|\ge\alpha\}$. We claim that
\begin{equation}\label{eq:ent-coupling}
 \{\tau_a^{\rm tr}\le T\}\subset\{\tau_a^{\rm cut}\le T\}
 \cup\{\varsigma^{\rm tr}\le T\}.
\end{equation}
If $|R(X_0,u)|\ge a$, both hitting times are zero. If $|s_0|\ge\alpha$, then
$\varsigma^{\rm tr}=0$. Otherwise, suppose that $\tau_a^{\rm tr}\le T$ and
$\varsigma^{\rm tr}>T$. Since $a<\alpha$ and the path is continuous, the
true process reaches $|r|=a$ before $|r|=\alpha$, and it does not reach
$|s|=\alpha$ before time $T$. The two processes therefore coincide on
$[0,\tau_a^{\rm tr}]$, and $\tau_a^{\rm cut}=\tau_a^{\rm tr}\le T$. This
proves \eqref{eq:ent-coupling}.

The law of $(J,u,X^{\rm cut})$ is $\mathsf Q_\alpha$. Since $\gamma\le
\gamma_0$, \cref{prop:ent-truncated} bounds the probability of the first
event in \eqref{eq:ent-coupling} by $C(N^{1-d/2}+N^{3-d})+e^{-cN}$. In the
same way, $\gamma\le\gamma_\alpha$ and the restricted bound of
\cref{lem:ent-boundary} imply that the second event has probability at most
$e^{-c_\alpha N}$.

\paragraph{Removing the restriction.}
Let now $J$ have its unrestricted Gaussian law, with the other variables as
before. For any event $\mathcal A$, we have
$\Pp(\mathcal A)\le\Pp(\mathcal A\mid\Omega_N)+\Pp(\Omega_N^c)$, and
$\Pp(\Omega_N^c)\le e^{-N}$. Under the unrestricted law and conditionally
on $u$, the field $H^{h,k,u}$ has the law $\E_{h,k,u}$ of
\cref{def:ent-conditioned}, which is the conditional law by
\cref{prop:ent-regression}, and it is independent of $X_0$ and $B$. By the measurable
construction of the solutions of \eqref{eq:ent-sde}, the probability of
$\{\tau_a^{\rm tr}\le T\}$ given $u$ is the left side of \eqref{eq:ent-main}
at the axis $u$; it is a measurable function of $u$ (compare
\cref{lem:red-kernels}). We have shown that
\begin{equation}\label{eq:ent-averaged}
 \int\E_{h,k,u}\Bigl[\Pp^H_{\nu_N}\Bigl(\sup_{0\le t\le e^{\gamma N}}
 |R(X_t,u)|\ge a\Bigr)\Bigr]\nu_N(\dd u)
 \le C\bigl(N^{1-d/2}+N^{3-d}\bigr)+e^{-cN}.
\end{equation}

\paragraph{Deterministic axes.}
Let $g\in O(N)$. The law of $H_N$ is invariant under $H\mapsto H\circ
g^{-1}$, and for the field $H\circ g^{-1}$ the quantities $H_Q(gu)$,
$K_Q(gu)$, and $G_{gu}$ equal $H_Q(u)$, $K_Q(u)$, and $gG_u$ for the field
$H$. By \eqref{eq:ent-deleted} and \eqref{eq:ent-conditioned}, the
conditioned field at $gu$ built from $H\circ g^{-1}$ is $H^{h,k,u}\circ
g^{-1}$. The law of $H^{h,k,gu}$ is therefore the image of the law of
$H^{h,k,u}$ under $H\mapsto H\circ g^{-1}$. If $X$ solves \eqref{eq:ent-sde} for a
Hamiltonian $\widetilde H$ with Brownian motion $B$, then $gX$ solves it for
$\widetilde H\circ g^{-1}$ with Brownian motion $gB$; this is the rotation
covariance of \cref{lem:red-kernels} at the level of paths. Together with
the invariance of $\nu_N$ and $R(gx,gu)=R(x,u)$, it shows that the
integrand in \eqref{eq:ent-averaged} takes the same value at $u$ and at
$gu$. As $O(N)$ acts transitively on $S_N$, the
integrand is constant, and \eqref{eq:ent-averaged} implies
\eqref{eq:ent-main} for every $u\in S_N$. In the pure case, the same argument applies with the
conditioning vector $(H_Q(u),G_u)$ and $k=ph$. Finite sums of exponentially small
terms are bounded by $e^{-cN}$ after decreasing $c$ and increasing $N_0$,
and this completes the proof.
\end{proof}
